\documentclass[UTF8,11pt,fontset=windows]{ctexart}
\usepackage[margin=2.35cm]{geometry}
\usepackage{amsmath,amssymb,unicode-math}
\usepackage{longtable,booktabs,array,multirow,calc,graphicx}
\usepackage[hidelinks]{hyperref}
\usepackage{bookmark,xurl,caption}
\captionsetup{labelformat=empty}
\setmainfont{Times New Roman}
\ctexset{section/format=\Large\heiti,subsection/format=\large\heiti,subsubsection/format=\normalsize\heiti}
\setcounter{secnumdepth}{-2}
\setlength{\emergencystretch}{3em}
\providecommand{\tightlist}{\setlength{\itemsep}{0pt}\setlength{\parskip}{0pt}}
\begin{document}



\hypertarget{paper-title}{%
\section{面向大语言模型推理的执行图优化研究综述}\label{paper-title}}

\hypertarget{abstract}{%
\subsection{摘要}\label{abstract}}

大语言模型推理需要反复执行结构相近的算子序列。随着算子计算不断加速，主机侧逐次提交、参数准备与同步逐渐成为影响推理时延的因素。ACLGraph和CUDA Graph通过预先组织任务及其依赖、重复提交执行图，减少逐算子调度开销。然而，真实服务中的批量变化、键值缓存增长、算子资源竞争和实例冷启动，使执行图优化超出了静态捕获与回放的范畴。本文以面向昇腾的ACLGraph应用为关注点，梳理执行图及相关编译、运行时研究，按优化对象划分为执行图构建与低开销提交、动态负载适配与图复用、并行执行与资源协调、图实例准备缓存与恢复四类，并进一步归纳十二个子类。本文比较各类方法改变的执行决策、依赖的运行条件及其引入的代价，并汇总代表研究的实际评测设置与结果边界，区分昇腾直接工作、CUDA Graph研究与可借鉴的其他系统方法。综合分析表明，图覆盖范围、动态适应能力、执行并行度与准备成本相互制约；单次回放加速不足以说明完整服务收益。未来研究需要结合负载分布和资源预算，联合考虑图构建、状态更新、缓存管理与恢复，并建立覆盖首次使用和稳态运行的统一评价方法。

\textbf{关键词}　大语言模型推理；ACLGraph；CUDA Graph；执行图；动态负载；任务调度；冷启动

\hypertarget{section-1}{%
\subsection{1 引言}\label{section-1}}

大语言模型的自回归推理包含提示词预填充和逐步解码两个主要阶段。预填充处理输入序列，解码则不断依据已有上下文生成新的token。尽管每轮输入内容与上下文长度有所变化，模型中的许多计算仍重复经过相似的算子序列。服务系统还需要在这些迭代之间调度请求、组织批量并维护键值缓存（key-value cache，KV缓存）。Orca通过迭代级调度适应请求的不同生成进度，PagedAttention则改善KV缓存的内存组织；这些工作同时说明，推理过程具有计算结构重复、运行状态持续变化的双重特点\protect\hyperlink{ref-1}{{[}1}, \protect\hyperlink{ref-2}{2{]}}。

在即时执行模式下，主机为算子调用准备参数并向加速器提交任务。单个算子耗时较短时，这些工作可能占据不可忽略的时间；即使设备上的计算已被优化，主机也未必能够及时提供后续任务。执行图将部分准备工作提前完成，把多个任务及其依赖组织成可重复提交的对象。CUDA Graph与ACLGraph均体现了这种思路，具体对象模型、可更新状态和资源限制则由各自运行时决定\protect\hyperlink{ref-3}{{[}3}, \protect\hyperlink{ref-4}{4{]}}。xLLM在昇腾上围绕部分图、动态维度和图内存池进行适配；GraCE则从编译器侧研究CUDA Graph的捕获兼容性、参数传递与收益选择\protect\hyperlink{ref-5}{{[}5}, \protect\hyperlink{ref-6}{6{]}}。它们共同表明，减少提交次数只是图执行设计的起点。

真实推理服务至少带来四组相互关联的问题。首先，注意力、自定义算子和运行时控制逻辑未必能进入同一张图，图边界决定了主机仍需参与多少工作。其次，批量、序列长度和KV布局发生变化后，原有图可能需要更新状态、选择其他变体或回退执行。再次，任务被统一提交并不意味着资源得到充分利用，独立算子之间仍可能存在不必要的串行与资源争用。最后，图捕获、编译、内存准备和通信初始化会形成启动成本，提前生成的图也未必能够在新进程中直接使用。这些问题分别涉及图的形成、复用、执行和存续，不能仅靠增加捕获范围解决。

已有LLM推理综述从模型、算法、系统和硬件等层次整理优化方法，或从实例与集群服务角度分析部署问题\protect\hyperlink{ref-7}{{[}7}, \protect\hyperlink{ref-8}{8{]}}。本文收窄研究对象，集中讨论运行时执行图相关的机制及其相互影响。研究范围既包含直接使用ACLGraph或CUDA Graph的工作，也包含有助于解释动态编译、细粒度调度和执行环境恢复的相邻方法。后者的纳入依据是其能够回答具体的执行图问题，而非论文中是否出现``graph''一词。

本文的组织目标有三点：建立具有明确判据的分类体系，使读者能够判断一项技术改变了什么；在同一问题下比较方法的执行过程与成立条件，避免逐篇罗列系统；以昇腾直接证据为基础，分析GPU研究中哪些机制具有参考价值、哪些仍依赖平台特定能力。第2节定义研究对象并介绍分类依据，第3至6节分别展开四类方法，第7节比较跨类特征、昇腾工程证据及实际评测结果，第8节归纳研究挑战，第9节总结全文。

\hypertarget{section-2}{%
\subsection{2 研究对象与分类体系}\label{section-2}}

\hypertarget{section-2-1}{%
\subsubsection{2.1 执行图及其基本过程}\label{section-2-1}}

本文以有向无环依赖图G=(V,E)作为基本分析模型。节点集合V表示设备计算、数据拷贝或同步等任务，边集合E表示任务之间必须满足的先后关系。节点还关联输入输出地址、运行参数和所需资源。只有这些信息与运行环境一致时，图才构成可执行的工作计划。本文所说的执行图不同于作为算法输入的图数据，也不同于表示候选token关系的投机解码树；编译器中间表示中的计算图与运行时图对象同样属于不同层次。编译器可以改写、融合或划分计算图，运行时再对生成的任务进行捕获和回放。PyTorch 2和Relax分别展示了程序提取与后端执行、不同层次中间表示组合的设计\protect\hyperlink{ref-9}{{[}9}, \protect\hyperlink{ref-10}{10{]}}。

图的建立通常通过显式描述任务依赖，或捕获已有流上的任务调用完成。准备阶段还可能包含算子编译、内存分配和实例化；回放阶段则复用已经准备的对象。CUDA文档将图定义、实例化和执行区分为不同阶段；昇腾文档以流任务捕获至运行实例、执行及销毁说明ACLGraph的基本接口过程\protect\hyperlink{ref-3}{{[}3}, \protect\hyperlink{ref-4}{4{]}}。这些阶段在不同框架中可能被封装到同一接口中，因此本文按功能区分它们，不要求平台具有完全相同的调用序列。

回放复用的是任务组织和已准备的运行状态，而非先前的计算结果。以解码为例，下一轮token值必须进入相应输入位置，KV缓存的长度和位置也必须更新。若图仍读取旧指针或旧长度，即使任务依赖本身没有变化，执行也可能不正确。相反，若新输入可以通过现有参数和缓冲区表达，就不一定需要重新捕获整张图。图复用的关键是确定哪些信息保持不变，以及哪些变化能够由已支持的更新机制表示。

图1给出贯穿全文的简化例子。考虑两个相邻解码轮次，活动请求数由3变为5，已有请求的KV长度增长，输入张量也可能位于新地址。该例只用于说明决策，不代表某个实测配置。在第一类方法中，系统决定注意力、归一化和前馈网络分别在哪些区域入图；第二类方法判断现有图能否容纳5个请求，并更新长度、块表或地址；第三类方法在依赖许可时安排计算与通信；第四类方法决定尚不存在的图何时捕获，以及进程重启后哪些准备结果可以恢复。由此，同一轮推理的延迟受到多个层次共同影响。

\begin{figure}
\centering
\includegraphics[width=\linewidth]{figures/execution-flow.png}
\caption{\textbf{图1　执行图复用的基本过程。实线表示一次执行必须经过的决策，虚线表示首次使用或恢复时发生的准备。批量3到5为概念示例，箭头不表示固定耗时或特定平台的API顺序。}}
\end{figure}

\hypertarget{section-2-2}{%
\subsubsection{2.2 分类依据与层次}\label{section-2-2}}

本文依据方法主要改变的执行决策划分四个大类。执行图构建与低开销提交决定哪些计算被组织成图、怎样建立可重复提交的计划，以及图模式是否值得采用；动态负载适配与图复用决定输入变化后如何选择执行方案并更新状态；并行执行与资源协调决定就绪工作如何推进、分配和共享资源；图实例准备缓存与恢复决定实例何时创建、如何重新可用，以及多个实例之间怎样减少重复准备。表1将这些问题细分为十二个小类，并列出对应机制。

第一类的三个决策依次是捕获合法性、计划形成方式和性能选择。兼容性改造回答一个区域能否进入图；显式构建或预执行捕获回答计划从哪里获得；候选实测回答已有合法方案中哪一种更有收益。按输入尺寸选取一个兼容图，不等于已经完成收益测量，因此将xLLM的形状驱动模式选择放入第二类，而以GraCE的候选实测代表性能选择。

第二类按需要改变的执行对象划分。判断时先检查执行路径是否变化，再检查既有内核与启动配置是否仍适用，最后检查同一配置内的数据、元数据和地址如何更新。拓扑或路径适配属于4.2节，规格专门化及变体选择属于4.1节，在所选执行结构内更新本轮状态属于4.3节。例如token数既可能触发换图，也可能仅作为参数写入；投机树变化也可能由掩码表达。归类依据是实现采用的处理机制，不是输入变量的名称。一个完整方案可以依次经历多个决策。

第三类按调度器作出的决策划分。并行计划与依赖约束决定哪些工作可以同时推进、哪些顺序必须保持；任务分配与负载均衡决定由哪个执行单元承担就绪工作，以及是否在运行时决定分派位置；资源共置与流水组织决定同时执行的工作占用多少资源、怎样交错使用计算与通信能力。任务粒度和单卡或多卡是横向属性。例如Mustard的跨设备依赖推进与节点内动态分配分别进入5.1和5.2节，MPK的细粒度任务可以预先派入指定worker队列，也可以等待依赖满足后再分派。更小的任务和更多设备本身都不是独立的优化目标。

第四类区分实例创建与缓存、跨进程恢复、以及多规格共享。6.1节处理实例尚未建立或缓存未命中的情形；6.2节比较新进程如何重建地址和内核绑定；6.3节讨论多个规格怎样共享结构和准备结果。Medusa的地址重绑定与Foundry的地址布局重现因而在同一小节比较，Foundry的模板去重则单独讨论。环境预热与通用进程恢复作为支撑机制说明其作用范围，不与图模板共享等同。

\begingroup\small
\begin{longtable}[]{@{}>{\raggedright\arraybackslash}p{(\linewidth - 4\tabcolsep) * \real{0.16}}>{\raggedright\arraybackslash}p{(\linewidth - 4\tabcolsep) * \real{0.26}}>{\raggedright\arraybackslash}p{(\linewidth - 4\tabcolsep) * \real{0.58}}@{}}
\caption{表1　执行图优化分类及代表机制}\tabularnewline
\toprule
\textbf{大类} & \textbf{小类与判据} & \textbf{代表机制及与执行图的关系} \\
\midrule
\endfirsthead
\toprule
\textbf{大类} & \textbf{小类与判据} & \textbf{代表机制及与执行图的关系} \\
\midrule
\endhead
执行图构建与低开销提交 & 捕获区域选择与兼容性改造：哪些任务能够入图 & 昇腾直接：xLLM部分图与算子参数化\protect\hyperlink{ref-5}{{[}5{]}}；CUDA直接：GraCE兼容性变换\protect\hyperlink{ref-6}{{[}6{]}}、Relax存储规划后卸载\protect\hyperlink{ref-10}{{[}10{]}} \\
& 执行计划构建与重复提交：怎样获得并复用计划 & CUDA直接：Nimble 2020预执行捕获\protect\hyperlink{ref-11}{{[}11{]}}、Taskflow显式图与调用捕获\protect\hyperlink{ref-12}{{[}12{]}} \\
& 图模式性能选择：合法方案中采用哪个版本 & CUDA直接：GraCE对无图、普通图和参数间接访问图实测选优\protect\hyperlink{ref-6}{{[}6{]}} \\
动态负载适配与图复用 & 形状专门化与图变体选择：怎样覆盖输入规格 & 昇腾直接：xLLM参数化与模式选择\protect\hyperlink{ref-5}{{[}5{]}}；CUDA直接：FlashInfer配置选图\protect\hyperlink{ref-13}{{[}13{]}}、Hydragen变长后缀兼容\protect\hyperlink{ref-14}{{[}14{]}}；编译支撑：Nimble 2021形状处理\protect\hyperlink{ref-15}{{[}15{]}} \\
& 动态路径与执行结构适配：路径变化怎样表达 & CUDA直接：Relax合格区域卸载\protect\hyperlink{ref-10}{{[}10{]}}、Sequoia预定树结构\protect\hyperlink{ref-16}{{[}16{]}}；编译支撑：Cortex线性化\protect\hyperlink{ref-17}{{[}17{]}}、ACRoBat动态批处理\protect\hyperlink{ref-18}{{[}18{]}} \\
& 回放状态与地址更新：已有方案如何读取本轮值 & 昇腾直接：xLLM固定偏移池\protect\hyperlink{ref-5}{{[}5{]}}；CUDA直接：FlashInfer计划工作区\protect\hyperlink{ref-13}{{[}13{]}}、GraCE参数间接访问\protect\hyperlink{ref-6}{{[}6{]}} \\
并行执行与资源协调 & 并行计划与依赖约束：哪些工作按何种顺序推进 & CUDA直接：Nimble流映射\protect\hyperlink{ref-11}{{[}11{]}}、Opara资源感知顺序\protect\hyperlink{ref-19}{{[}19{]}}、AutoGraph并行图改写\protect\hyperlink{ref-20}{{[}20{]}}、Mustard跨设备依赖\protect\hyperlink{ref-21}{{[}21{]}}；调度支撑：IOS阶段选择\protect\hyperlink{ref-22}{{[}22{]}} \\
& 任务分配与负载均衡：由哪个执行单元执行 & CUDA直接：Mustard设备端子图分派\protect\hyperlink{ref-21}{{[}21{]}}；执行模型支撑：Rammer静态rProgram\protect\hyperlink{ref-23}{{[}23{]}}、MPK事件驱动任务分配\protect\hyperlink{ref-24}{{[}24{]}} \\
& 资源共置与流水组织：并发工作怎样共享资源 & 调度支撑：NanoFlow子批次流水\protect\hyperlink{ref-25}{{[}25{]}}、POD-Attention线程块共置\protect\hyperlink{ref-26}{{[}26{]}}；昇腾系统支撑：xLLM Cube与Vector分配\protect\hyperlink{ref-5}{{[}5{]}} \\
图实例准备缓存与恢复 & 图实例创建与缓存管理：何时生成和保留实例 & 昇腾直接：xLLM多缓存部分图\protect\hyperlink{ref-5}{{[}5{]}}；CUDA直接：Hybrid按需捕获与滚动缓存\protect\hyperlink{ref-27}{{[}27{]}} \\
& 跨进程图恢复与地址策略：怎样恢复有效绑定 & CUDA直接：Medusa分配索引重绑定\protect\hyperlink{ref-28}{{[}28{]}}、Foundry地址布局和内核环境重建\protect\hyperlink{ref-29}{{[}29{]}} \\
& 图模板共享与重复准备消除：多个规格怎样共享准备 & CUDA直接：Foundry同拓扑模板与参数集\protect\hyperlink{ref-29}{{[}29{]}}；环境支撑：DEEPSERVE预热引擎\protect\hyperlink{ref-30}{{[}30{]}}、PhoenixOS上下文准备\protect\hyperlink{ref-31}{{[}31{]}} \\
\bottomrule
\end{longtable}
\endgroup

表1中的关系标记描述具体机制，不是论文质量等级。一项机制以主要改变的决策确定归属；一个系统可以包含多项机制，因而在不同位置出现。xLLM的捕获与状态机制属于ACLGraph直接证据，其Cube与Vector协调则属于昇腾资源优化。平台、任务粒度、决策时机和验证负载作为横向属性比较，既保留研究之间的交叉关系，也避免用系统名称代替分类依据。

\hypertarget{section-2-3}{%
\subsubsection{2.3 文献范围与比较口径}\label{section-2-3}}

本文以截至2026年11月23日整理的论文全文和版本固定的官方材料为基础，覆盖运行时执行图的捕获、回放、适配、调度与恢复。文献选择采用问题导向方式：以直接研究ACLGraph或CUDA Graph机制的论文为核心，结合其相关工作检查相邻编译、任务运行时和恢复技术，再以官方接口及公开实现核实昇腾工程条件。检索和筛选对象包括相关会议论文、作者预印本、官方文档与代码仓库，本文不据此声称已完成全领域穷尽检索。

纳入材料分为三种用途。核心工作必须说明如何改变图的构造、回放、状态、调度或准备过程。支撑工作需要提供与核心问题直接对应的替代机制或成立条件，例如Nimble 2021的动态形状计算、MPK的设备端任务分配、通用GPU检查点的恢复边界；讨论止于这些关系，不将其整体性能归为图优化收益。背景工作用于交代输入变化和服务目标，例如Orca的请求迭代与PagedAttention的KV映射。仅共享``graph''名称、或只优化模型压缩而未建立上述关系的工作，不作为执行图方法代表。

支撑文献的加入应补充新的比较维度。比如，S-LoRA中的异构适配器与分页存储、Punica的分段低秩计算，以及Sarathi-Serve的分块预填充说明，输入变化还涉及适配器组成和每轮预填充/解码混合比例。本文在评测范围讨论中纳入这些负载，未据此将请求调度或LoRA计算列为图捕获机制\protect\hyperlink{ref-32}{{[}32}, \protect\hyperlink{ref-33}{33}, \protect\hyperlink{ref-34}{34{]}}。论文与官方实现也承担不同证据角色：论文提供其评测范围内的性能结果，文档与源码说明某版本中实现了哪些动作，二者不能互相替代。

为统一收益讨论，考虑固定模型和配置重复执行N轮的简化场景。记\emph{T\textsubscript{p}}为相对于即时执行新增的一次性准备成本，\emph{t\textsubscript{e}}与\emph{t\textsubscript{g}}分别为即时路径和图路径的平均端到端单轮耗时，则图路径具有总时间收益的条件为：

\hypertarget{equation-1}{}
\begin{equation}T_p+N\times t_g<N\times t_e\tag{1}\end{equation}

当\emph{t\textsubscript{e}}\textgreater{}\emph{t\textsubscript{g}}时，式（1）可写为N\textgreater{}\emph{T\textsubscript{p}}/(\emph{t\textsubscript{e}}−\emph{t\textsubscript{g}})。这里\emph{t\textsubscript{g}}包含图选择、必要的数据与参数更新、提交、设备执行及同步的综合时间；不能只取设备内核耗时。该式用于解释摊销关系，不对任意系统给出实测预测。动态负载下，图未命中、后台准备、排队和资源干扰会改变各轮成本，需要结合请求轨迹记录。第7节因此同时列出原论文实际评测的条件和结果，保留收益不明显或出现退化的情形，并区分系统整体收益与图机制单项收益。

比较一项机制时，本文依次考察其省去了什么工作、增加了什么工作，以及何种输入与环境使两者的差额有利。固定缓冲区减少了地址变化，却增加复制；图专门化减少填充，却增加准备和保留资源；运行时分配改善不均衡，却需要队列与同步；环境物化迁移了在线准备，却要求额外存档和兼容性检查。不同硬件和基线中的加速比用于描述各自实验，不直接排列成统一性能名次。

\hypertarget{section-3}{%
\subsection{3 执行图构建与低开销提交}\label{section-3}}

一次解码迭代包含输入准备、模型计算和结果处理。即使模型相同，主机仍可能反复解释算子、安排执行流、准备参数并提交设备任务。执行图将其中可复用的工作提前完成，使后续迭代通过较少的主机调用启动计算。这一过程包含三个不同决策：哪些区域能够进入图，怎样获得可重复提交的执行计划，以及在若干正确的执行模式中选择哪一种。前两个决策建立可用路径，第三个决策检验其性能收益；扩大捕获范围并不自动意味着每轮推理更快。

\hypertarget{section-3-1}{%
\subsubsection{3.1 捕获区域选择与兼容性改造}\label{section-3-1}}

模型级计算图描述张量之间的运算关系，设备执行图则需要记录具体内核、参数、内存访问及同步关系。前者能够被编译器追踪，并不保证后者可以捕获。例如，算子内部可能依据输入形状执行主机计算、临时构造 CPU 张量，或在运行中申请工作空间。面对此类障碍，可以在现有实现中提取可捕获区域，也可以改变算子和存储组织，使原本不兼容的区域进入图。两种方法分别调整捕获边界和改造内部实现，使所选区域满足捕获约束，能够作用于同一模型的不同部分。

xLLM 提供了昇腾上的直接实例。其 Adaptive Graph Mode 基于 ACLGraph，将动态形状较简单的前馈网络（FFN）、层归一化（LayerNorm）等模块组织为部分图，而在所描述方案中，将具有多个动态维度的注意力模块保留在即时执行路径。这样可以先复用稳定区域，无需等待全部算子获得相同的动态支持。对于 PageAttention、AllReduce 等自定义算子，xLLM 又识别其原本在主机上即时计算的形状信息，将维度和循环次数改为可由图输入导出的内核参数，从而扩大捕获范围。这里的注意力图外执行是特定实现的边界；自定义算子改造本身已经说明，这一边界可以随实现改变\protect\hyperlink{ref-5}{{[}5{]}}。

GraCE 处理的主要障碍位于通用编译框架内部。在其研究的 PyTorch 2 路径中，CPU 标量、部分主机与设备间复制及相关设备放置信息，会使整个 FX 图失去 CUDA Graph 捕获资格。GraCE 沿中间表示与源程序的映射定位原因，将特定标量构造为设备张量，或把可以提前完成的复制移出捕获路径，再重新生成中间表示。相较于模块级部分图策略，这一路径优先修复使已有编译区域失去资格的操作和元数据。其成立条件是编译器能够识别相应行为并安全转换；论文列举的兼容性变换不能推广为任意主机逻辑均可入图\protect\hyperlink{ref-6}{{[}6{]}}。

Relax 则从存储准备入手改变可捕获范围。它保留符号形状关系，利用尺寸等价和活跃期分析规划存储；动态维度具有已知上界时，可以提前分配足够空间。对于张量程序内部申请的全局工作空间，Relax 联合改写内核与调用方，将分配提升到图层，再纳入统一规划。在这些变换之后，编译过程才提取满足 CUDA Graph 条件的子图，并插入捕获和回放调用。因此，Relax 解决的是如何让存储条件可提前确定，不能仅由``支持动态形状''推导出任意动态程序都可以捕获\protect\hyperlink{ref-10}{{[}10{]}}。

三项工作可以在相同障碍下比较。若无法捕获的是含主机逻辑的模块，区域选择能够保留其他部分的收益；若障碍来自可识别的标量、复制或参数形式，代码变换有机会消除图外边界；若障碍是中间存储和工作空间，则需要先改写分配位置。选择还应检查新增代价：部分图增加区域切换，兼容性改造引入编译约束，预分配扩大常驻存储。因而本节的判据不是单独最大化入图算子比例，而是在明确障碍及改造成本后，选择能够减少重复提交且适配成本可控的区域。上述选择是对机制的归纳，具体收益仍由第3.3节的测量决定\protect\hyperlink{ref-5}{{[}5}, \protect\hyperlink{ref-6}{6}, \protect\hyperlink{ref-10}{10{]}}。

\hypertarget{section-3-2}{%
\subsubsection{3.2 执行计划构建与重复提交}\label{section-3-2}}

捕获区域确定后，还需获得可以反复启动的任务组织。关键差别在于计划从哪里来：从一次已有执行中记录任务轨迹，还是直接描述任务及依赖。前者能够承接框架和库已经生成的调用，后者提供显式控制；两者均需保留回放所依赖的内存和设备环境。本节讨论计划的形成与提交方式，图何时创建、保留多久及未命中后如何处理属于第6节的实例管理。

Nimble（Kwon 等，2020）以执行相同计算集合的静态神经网络为对象。它首先分析计算图，确定算子与执行流的映射并插入同步，再用虚拟输入预执行模型，截获设备任务和内存申请。记录内容包括内核及参数、提交次序和所属流，预执行期间使用的设备内存也被保留。其实现使用流捕获建立 CUDA Graph，随后通过图启动接口重复提交，绕过原框架逐轮安排任务与分配内存的工作。计划能够复用的原因在于，运行时重新执行的是已经确定的设备任务，而非再次发现这些任务\protect\hyperlink{ref-11}{{[}11{]}}。

记录执行轨迹也意味着计划带有适用条件：任务集合、依赖及存储安排需要继续成立。当输入变化导致路径或资源需求超出原计划时，预执行本身不会自动建立新的兼容关系。图1中请求数由3增至5，就需要进一步检查既有计划是否覆盖新的规模。这里的 Nimble 是2020年的 GPU 任务调度系统；第4节讨论的 Nimble（Shen 等，2021）通过动态编译和虚拟机处理变化的程序，两者属于不同研究，不能把后者的动态能力归给前者\protect\hyperlink{ref-11}{{[}11}, \protect\hyperlink{ref-15}{15{]}}。

Taskflow 提供显式描述与捕获已有调用两种入口。其 cudaFlow 允许程序给出内核、复制操作及依赖，并将 GPU 子图嵌入更大的 CPU---GPU 任务图。能够获得内核参数时，应用可直接构建 CUDA Graph；对于 cuDNN、cuBLAS 等主要暴露流接口、内部内核不可见的库，cudaFlow capturer 捕获调用产生的设备任务。两种方式因而对应不同的信息可见性：显式构图需要掌握任务描述，流捕获沿用库已经实现的发射路径。应用仍负责指针、设备内存放置及相关环境，编程接口并不会消除这些正确性条件\protect\hyperlink{ref-12}{{[}12{]}}。

在相同的重复任务序列下，Nimble 的重点是从模型预执行自动获得计划，Taskflow 的重点是把显式任务描述、库调用和异构程序组合起来。Taskflow 的条件任务可以在 CPU 与 GPU 子图边界组织循环和分支，但这不意味着所有控制决策都位于同一个 CUDA Graph 内。Relax 的首次捕获、后续回放也属于这一准备与执行分离模式，其计划来自编译器筛选的合格子图。若主要目标是承接已有模型调用，自动记录能够减少手工描述；若需要控制 CPU---GPU 任务组合，显式接口提供更直接的组织方式。无论选择哪条路径，动态状态均需另行适配，计划自动化程度不能替代复用范围的检查\protect\hyperlink{ref-10}{{[}10}, \protect\hyperlink{ref-11}{11}, \protect\hyperlink{ref-12}{12{]}}。

\hypertarget{section-3-3}{%
\subsubsection{3.3 图模式性能选择}\label{section-3-3}}

当同一区域已有若干正确的执行路径时，问题才转为性能选择。逐内核提交能够被合并，但图回放之前可能需要复制输入、更新指针，并承担图对象管理等运行时工作。若主机发射开销本来较小，或数据准备较重，新增代价就可能抵消提交节省。本节将``根据输入选择兼容实例''与``比较正确候选的耗时''分开；前者在第4.1节讨论，后者以 GraCE 的明确实测机制为核心。

GraCE 在编译阶段为同一模块考虑至多三种候选：无 CUDA Graph、普通 CUDA Graph，以及使用参数间接访问的 CUDA Graph。其分析器在目标设备上运行可用版本、测量执行时间，并缓存实测最快的版本供后续调用使用。三者构成一个有共同任务和测量对象的对照：无图保留逐次发射而避免图相关准备；普通图合并发射但需要维护固定参数位置；间接访问图进一步把部分数据复制换为指针更新。选择对象是某个编译模块的执行方式，并非不同模型或不同硬件上的加速比\protect\hyperlink{ref-6}{{[}6{]}}。

保留普通图候选是必要的。GraCE 的指针更新来自主机，经过主机到设备的传输路径；被替代的数据复制则可能发生在设备内部。处理许多小对象时，指针占用字节较少，并不保证传输及调用开销较小。对于可修改的 JIT 内核，间接访问可以直接写入内核；对于不可修改的库内核，则需要前置内核更新节点参数。后者还增加设备任务，进一步改变两种图模式的成本关系。目标硬件的传输能力、对象数量以及内核可改写性，因此都进入实际选择，而不能仅凭减少了复制字节数决定采用间接访问\protect\hyperlink{ref-6}{{[}6{]}}。

这一选择发生在编译慢路径，测量和准备成本仍需要后续调用摊销。第2节式（1）区分了新增准备时间与平均单轮耗时：即使选出的候选在稳态更快，短时任务也未必获得总时间收益。GraCE 给出的是目标硬件上局部候选的实测选择，不能单独覆盖未来输入分布、图缓存淘汰和进程重启的全部成本。将其与实例管理组合时，还需说明候选测量针对什么配置、结果何时复用，以及配置改变后是否重新选择。这里的重新选择条件属于系统需要明确的管理规则，不能补写为 GraCE 已经实现的在线策略\protect\hyperlink{ref-6}{{[}6{]}}。

据此，图模式性能选择应当在捕获合法性检查之后进行，并计入候选各自的输入准备、提交及设备执行时间。若数据复制是主要代价，可考察间接访问；若模块任务较短、图外处理占比高，则应保留无图路径参与比较；若准备结果复用次数有限，还需结合完整使用周期决定是否投入候选搜索。xLLM 的形状驱动模式选择具有不同决策依据，本文在第4.1节将其作为动态适配机制讨论，不据其模式切换推导论文未给出的性能搜索算法\protect\hyperlink{ref-5}{{[}5}, \protect\hyperlink{ref-6}{6{]}}。

表2按具体决策组织机制，使同一问题下的方法能够直接对照。捕获改造、计划生成和性能选优可以先后作用于同一解码区域，表中的三组方法因此主要是互补关系。

\begingroup\small
\begin{longtable}[]{@{}>{\raggedright\arraybackslash}p{(\linewidth - 6\tabcolsep) * \real{0.19}}>{\raggedright\arraybackslash}p{(\linewidth - 6\tabcolsep) * \real{0.27}}>{\raggedright\arraybackslash}p{(\linewidth - 6\tabcolsep) * \real{0.27}}>{\raggedright\arraybackslash}p{(\linewidth - 6\tabcolsep) * \real{0.27}}@{}}
\caption{表2　执行图构建与提交机制比较}\tabularnewline
\toprule
\textbf{决策与工作} & \textbf{计划或障碍来源} & \textbf{处理机制} & \textbf{主要条件与代价} \\
\midrule
\endfirsthead
\toprule
\textbf{决策与工作} & \textbf{计划或障碍来源} & \textbf{处理机制} & \textbf{主要条件与代价} \\
\midrule
\endhead
区域选择：xLLM（昇腾）\protect\hyperlink{ref-5}{{[}5{]}} & 动态算子及主机形状逻辑 & 部分图；自定义算子参数化 & 算子须支持；保留图外切换 \\
兼容改造：GraCE\protect\hyperlink{ref-6}{{[}6{]}} & 标量、复制和设备元数据 & 定位障碍并改写程序 & 需要安全的编译变换 \\
存储准备：Relax\protect\hyperlink{ref-10}{{[}10{]}} & 动态分配及内核工作空间 & 上界规划；工作空间提升 & 需要尺寸关系；预留容量 \\
轨迹记录：Nimble（2020）\protect\hyperlink{ref-11}{{[}11{]}} & 模型预执行产生的任务 & 记录多流任务并保留内存 & 计算集合和存储条件稳定 \\
任务描述：Taskflow\protect\hyperlink{ref-12}{{[}12{]}} & 显式依赖或库的流调用 & 显式构图／流捕获 & 应用维护内存与上下文 \\
性能选优：GraCE\protect\hyperlink{ref-6}{{[}6{]}} & 同一模块的正确候选 & 实测无图、普通图、间接访问图 & 测量与准备需要摊销 \\
\bottomrule
\end{longtable}
\endgroup

\hypertarget{section-4}{%
\subsection{4 动态负载适配与图复用}\label{section-4}}

执行图复用任务组织，推理请求却持续改变执行条件。连续批处理允许请求逐轮加入和退出，序列长度随生成增长，投机推理还会改变候选 token 及其可见关系。Orca 的迭代级调度揭示了这种动态性：模型权重和算子类型不变，并不意味着每轮输入组成相同\protect\hyperlink{ref-1}{{[}1{]}}。

本章按适配动作区分三个问题：现有配置是否仍适合输入，动态路径能否由可复用结构表达，以及选定结构怎样读取本轮状态。它们对应形状专门化与变体选择、路径与结构适配、状态与地址更新。一项输入变化可能依次触发多个动作。例如图1中批量由3增至5，系统可能先换图，再更新长度与地址；若同一参数化图已经覆盖两种规模，则只需后续更新。因而是否改变图，应依据内核、任务配置和更新接口判断，不能只看输入形状是否变化。

\hypertarget{section-4-1}{%
\subsubsection{4.1 形状专门化与图变体选择}\label{section-4-1}}

形状适配在专门化效率与覆盖范围之间取舍。固定尺寸有利于选择内核、组织启动配置和规划存储，但为每个输入准备独立执行计划会增加编译、捕获及保留成本。相反，单一计划覆盖更多形状，需要内核接受动态参数，或承担填充和通用代码的代价。因此，比较对象应是同一形状变化由多大范围的计划覆盖，以及覆盖所付出的运行和准备成本。

xLLM 在昇腾上结合动态维度参数化与多缓存部分图。对于动态维度较简单的 FFN、LayerNorm 等模块，将 token 数等变化量作为图输入，用于内存需求计算和内核配置；对于所述实现中动态支持更复杂的注意力模块，保留即时执行。系统依据当前形状在即时执行、完整图和参数化部分图之间选择，并缓存部分图供后续使用。它强调的是输入与模式之间的适配关系，报告未给出类似 GraCE 的同模块候选实测算法。可覆盖哪些形状仍取决于算子支持，不能把参数化理解为所有维度都能在任意图实例中改变\protect\hyperlink{ref-5}{{[}5{]}}。

FlashInfer 将配置专门化与逐轮规划分开。它依据平均查询长度和 KV 组合格式等配置编译内核，捕获相应的 CUDA Graph，由服务框架根据当前 KV 配置选择图实例。选定图后，具体序列长度由每轮生成的计划处理。这使图实例的选择尺度与运行时状态的变化尺度不同：改变 KV 格式或适宜的内核配置可能需要换图，长度增长却可以在选定图中通过计划更新处理。与 xLLM 相比，两者均扩大重复使用机会，但前者突出注意力配置及规划分离，后者突出部分图与算子动态维度接口，不能仅按缓存图数量判断优劣\protect\hyperlink{ref-5}{{[}5}, \protect\hyperlink{ref-13}{13{]}}。

Hydragen 提供了动态长度保持图兼容的具体内核实例。其主要方法将共享前缀与独有后缀的注意力拆开，跨序列批量计算共享前缀，再合并结果。在实现中，前缀使用 flash-attn，后缀使用 xformers 的 Triton 内核；论文第3.5节明确指出，后者允许后缀 KV 长度随解码步增长，同时满足 CUDA Graph 的约束。这一事实说明，某个长度维度变化不必导致每轮重新捕获，但论文没有据此给出通用图变体缓存算法。其共享前缀计算收益也不能归因于图兼容本身；所述原型仍要求用户指定输入中的共享关系\protect\hyperlink{ref-14}{{[}14{]}}。

更广的形状覆盖需要编译器在不知道具体尺寸时保留可用信息。Nimble（Shen 等，2021）以 Any 表示编译时未知的维度，在运行时通过形状函数计算输出尺寸并完成必要检查。针对符号维度的分块循环，它按尺寸对分块因子的余数生成不同内核，再依据实际形状选择对应实现，以便化简循环边界检查；候选也可包含第三方库实现。由此，专门化不必为每个完整尺寸分别编译，而可以围绕影响代码效率的条件组织。其调优还在若干形状上比较候选的平均表现；减少候选能够控制代码规模，但会损失部分专门化机会。该工作生成的是动态程序与内核，不是 CUDA Graph 缓存器\protect\hyperlink{ref-15}{{[}15{]}}。

Nimble、DISC 与 Relax 对动态执行成本的处理位置不同。Nimble 将形状计算、分配和调用组织为虚拟机指令；DISC 则把形状推导、缓冲区管理、启动配置及库调用等主机逻辑和设备计算一起编译，减少解释执行，并利用维度相等、元素总数相等等约束辅助融合。DISC 针对动态尺寸、静态秩的张量，其动态支持不能直接解释为已有运行时图的回放能力。Relax 的跨层表示进一步保留符号表达式，利用尺寸等价和活跃期分析复用存储，并在上界已知时形成预分配计划，随后提取合格的 CUDA Graph 子图。三者分别提供动态程序执行、主机流程代码生成和面向图卸载的存储准备，其关系不是``支持动态形状''一个标签所能概括\protect\hyperlink{ref-10}{{[}10}, \protect\hyperlink{ref-15}{15}, \protect\hyperlink{ref-35}{35{]}}。

CoRa 从表示层减少覆盖不同长度所需的无效计算。它利用变长张量表示和局部填充，避免把所有不等长输入统一扩展到同一最大尺寸。这为比较填充与专门化提供了参照，但 CoRa 直接优化的是张量程序；其结果是否能进入可回放图，还需核实后续内存与捕获条件\protect\hyperlink{ref-36}{{[}36{]}}。

回到图1，若系统仅准备容量为4和8的图，批量3可以匹配前者，批量5则选择后者并处理填充；若同一参数化图已覆盖两种规模，便可能继续使用。容量仅为概念示例，不对应上述论文的具体分桶设置。选择规则应共同考虑形状覆盖、填充工作量、内核效率、变体数量和准备成本：范围稳定且专门化收益明显时，可以保留少量针对性变体；细小变化可由既有内核参数或计划表达时，应先检查能否继续复用。输入怎样匹配已有图属于本节，缺失图何时创建与淘汰则属于第6节。

\hypertarget{section-4-2}{%
\subsubsection{4.2 动态路径与执行结构适配}\label{section-4-2}}

形状变化未必改变执行路径，而条件分支、递归和投机验证可能改变实际需要完成的工作。适配首先要确定路径何时可知：输入结构已知时，可以提前形成计划；路径依赖中间张量值时，则必须在相应结果产生后继续决策。稳定区域分段执行、显式动态程序和结构数据化代表不同处理方式。比较它们时，需要同时检查保留了多少动态决策，以及准备和跨区域交互增加了什么工作。

区域化编译在动态决策之间提取稳定计算。Relax 用无控制流、无副作用的数据流区域组织优化，再从满足存储及捕获条件的程序中提取 CUDA Graph 子图。PyTorch 2 的 Dynamo 在遇到无法安全追踪的操作时形成图中断，编译已获得区域，并在之后继续追踪；GraCE 在相应流程内改造部分捕获障碍。这些方法允许控制逻辑继续留在外层，不要求整段动态程序成为一张固定执行图。代价是区域较碎时仍需跨越较多主机边界，因而不能只报告区域内部的提交节省\protect\hyperlink{ref-6}{{[}6}, \protect\hyperlink{ref-9}{9}, \protect\hyperlink{ref-10}{10{]}}。

Nimble（2021）采用显式动态程序承载控制，而非对所有路径预先展开。编译产物包含平台无关字节码和平台相关内核，虚拟机按指令执行控制逻辑、分配和内核调用。形状函数进一步区分仅依赖输入形状、依赖输入数值，以及先估计上界再取得实际形状的情况。若某个形状函数必须读取融合区域内的中间结果，论文会限制相应融合，避免在结果产生之前要求其形状已知。这揭示了动态适配的具体边界：提前编译控制代码能够减少框架组织开销，却不能消除由真实数据依赖决定的等待，也不等于把这些控制指令捕获进 CUDA Graph\protect\hyperlink{ref-15}{{[}15{]}}。

Cortex 与 ACRoBat 进一步展示了路径可知时机对方法的影响。Cortex 面向控制流依赖已知输入结构的递归模型，先在 CPU 上线性化结构，再把递归降低为循环，实施批处理、分支专门化和展开。ACRoBat 面向包含张量值相关控制流的动态程序，通过混合静态与动态分析发掘批处理机会，并利用编译后的控制代码、执行纤程及融合 gather 降低批次组织成本。前者能够提前整理执行次序，后者需要在运行中发现和组合工作。二者属于相关编译方法，提供结构转换与组织策略；自动批处理或递归降低本身不能证明转换结果可被执行图捕获\protect\hyperlink{ref-17}{{[}17}, \protect\hyperlink{ref-18}{18{]}}。

投机推理还允许把部分结构变化转为数据变化。SpecInfer 用树拓扑组织 KV，并通过拓扑感知因果掩码，在单个注意力内核中处理整棵候选树。EAGLE-2 根据上下文调整草稿树，随后将选定 token 展平成一维输入，并构造仅访问祖先节点的掩码。候选树的边由此可以影响索引和掩码，而不必与设备任务图的边一一对应。执行结构能否保持不变，仍取决于输入容量、启动配置及相应内核的兼容范围；这两项工作的树处理机制不能直接充当通用图复用证明\protect\hyperlink{ref-37}{{[}37}, \protect\hyperlink{ref-38}{38{]}}。

Sequoia 提供了更明确的直接联系：它通过动态规划预先选择投机树结构，并利用结构预先确定这一条件使用 PyTorch CUDA Graphs 降低启动开销。与根据上下文持续改变草稿树的方法相比，预定结构更容易形成可重复执行的任务组织，但需要接受其对结构适应性的限制；索引和掩码方案则需要为变化关系准备辅助数据。这不是固定树必然优于动态树的性能结论，而是两种方式对图复用施加了不同条件\protect\hyperlink{ref-16}{{[}16}, \protect\hyperlink{ref-38}{38{]}}。

因此，选择路径适配方法时应先确认动态性来自程序分支还是输入关系。若路径运行前已知，可以准备计划或线性化结构；若依赖尚未产生的结果，则保留动态控制，并尽量复用决策之间的稳定计算；若关系可以由固定内核消费的索引、掩码表达，则把结构变化下移为数据更新。最后一条路线必须继续检查第4.3节中的容量、地址和同步条件，不能仅凭数据已经展平便认定图可以安全复用。

\hypertarget{section-4-3}{%
\subsubsection{4.3 回放状态与地址更新}\label{section-4-3}}

选定可兼容的任务组织后，回放仍需读到本轮数据。这里区分输入与参数值、描述工作组织的元数据，以及这些对象所在的地址。图1中的 KV 长度增长、活动请求变化和输入张量迁移，可能同时修改三者。比较状态更新方法时，需要先固定问题：在任务及启动配置仍适用的条件下，哪些内容必须变化，哪些位置继续保持稳定；否则减少参数传输与减少完整输入复制就容易被误认为同一种收益。

xLLM 通过固定偏移的 HBM 池承接外部地址变化。初始化时预留连续空间，以相对池基址的固定偏移管理图内输入、输出和中间张量。回放前把用户输入复制到预定位置，执行后再把输出复制到调用方缓冲区。外部张量因而可以更换地址，图内仍使用已知布局。对具有主机侧即时形状计算的 PageAttention、AllReduce 等自定义算子，则通过改造内核参数传递动态维度或循环次数。前者维持数据位置，后者使动态数值可被内核使用；二者解决不同问题，共同代价包括固定容量和回放前后的搬运\protect\hyperlink{ref-5}{{[}5{]}}。

FlashInfer 把逐轮变化集中到计划元数据。CPU 侧 plan 依据序列长度生成线程块工作队列及部分结果到最终输出的映射，将计划异步复制到 GPU 工作区，run 再读取计划完成计算。相同长度规格的多层注意力可以复用计划，摊薄其准备成本。相关注意力及归约执行使用固定的网格大小，计划和中间结果在工作区中使用固定偏移，使传入内核的位置保持稳定。变化的是缓冲区中的工作描述，CPU 规划与计划传输则仍位于图捕获之外，必须计入每轮成本\protect\hyperlink{ref-13}{{[}13{]}}。

GraCE 针对的是图外输入地址改变所引起的数据复制。它捕获稳定的指针存储位置，每轮把当前数据地址写入该位置，内核通过间接访问取得本轮数据。对于可修改的 JIT 内核，编译器直接改写参数访问；对不可修改的库内核，则在图中加入前置内核，利用设备侧节点参数更新能力填入实际地址。指针复制与后续回放放在同一 CUDA 流中，通过流内顺序维持依赖。该机制减少的是把外部输入搬到固定数据区的需要，代价则转为主机到设备的指针传输、解引用，或额外的前置内核；是否采用仍需第3.3节的候选测量\protect\hyperlink{ref-6}{{[}6{]}}。

在相同条件下对照三者，可以辨认互补和替代关系。设所选图仍能处理当前批量，输入内容每轮变化且外部地址可变，KV 长度也在增长。xLLM 将输入值送到固定位置，解决外部地址不稳定；GraCE 允许内核间接访问外部数据，在机制层面提供了替代完整输入复制的思路，但它并不生成注意力的工作分配计划。FlashInfer 则解决长度变化下的任务组织，把队列和归约映射写入固定工作区；它的计划更新不能独自替代全部输入地址管理。因而不能按更新字节数给三项工作直接排序，应分别记录完整输入的搬运、计划生成及传输、指针更新和新增设备任务，再检查这些机制能否在同一实现中组合\protect\hyperlink{ref-5}{{[}5}, \protect\hyperlink{ref-6}{6}, \protect\hyperlink{ref-13}{13{]}}。

\begin{figure}
\centering
\includegraphics[width=\linewidth]{figures/state-update.png}
\caption{\textbf{图2　回放状态更新的三种路径。固定缓冲区复制以稳定的数据位置承接外部输入；计划元数据更新以稳定工作区承接变化的任务组织；指针间接访问以稳定的指针入口指向本轮数据。依据 xLLM、FlashInfer 与 GraCE 归纳，示意图不表示统一实现、传输量比例或实测耗时}\protect\hyperlink{ref-5}{{[}5}, \protect\hyperlink{ref-6}{6}, \protect\hyperlink{ref-13}{13{]}}\textbf{。}}
\end{figure}

这些路径还需要共同满足状态生命周期条件。固定地址不能防止下一轮覆盖尚未消费的数据，指针有效也不保证其指向的内容已准备完成；共享工作区若被多个回放使用，同样需要安排更新、读取和回收顺序。GraCE 给出了同流指针更新与回放的具体顺序方案，而更一般的多流、多实例组合仍须按实际实现检查。PagedAttention 的块映射与 vAttention 的物理页映射解释了 KV 状态如何改变，本节关心的是运行中的图怎样消费这些状态，并不把通用 KV 管理等同于图参数更新\protect\hyperlink{ref-2}{{[}2}, \protect\hyperlink{ref-6}{6}, \protect\hyperlink{ref-39}{39{]}}。

表3按适配动作归纳可复用对象和新增工作。若主要变化是外部数据地址，固定区复制与间接访问提供直接对照；若变化是工作划分，应先准备可由稳定内核消费的元数据；若变化超出启动配置或路径的兼容范围，则返回第4.1或4.2节重新选择实例或组织结构。这样，图复用是否有益取决于适配后的完整迭代成本，而非单独的图启动时间或缓存命中率。

\begingroup\small
\begin{longtable}[]{@{}>{\raggedright\arraybackslash}p{(\linewidth - 6\tabcolsep) * \real{0.16}}>{\raggedright\arraybackslash}p{(\linewidth - 6\tabcolsep) * \real{0.28}}>{\raggedright\arraybackslash}p{(\linewidth - 6\tabcolsep) * \real{0.28}}>{\raggedright\arraybackslash}p{(\linewidth - 6\tabcolsep) * \real{0.28}}@{}}
\caption{表3　动态适配机制比较}\tabularnewline
\toprule
\textbf{适配动作与工作} & \textbf{复用对象} & \textbf{运行时变化} & \textbf{新增工作或条件} \\
\midrule
\endfirsthead
\toprule
\textbf{适配动作与工作} & \textbf{复用对象} & \textbf{运行时变化} & \textbf{新增工作或条件} \\
\midrule
\endhead
形状匹配：xLLM（ACLGraph）\protect\hyperlink{ref-5}{{[}5{]}} & 参数化部分图 & 输入规模与模式 & 算子支持；多图准备 \\
配置选择：FlashInfer（CUDA Graph）\protect\hyperlink{ref-13}{{[}13{]}} & 配置专门化图 & KV配置与逐轮计划 & 多变体；CPU规划 \\
长度适配：Hydragen（CUDA Graph兼容）\protect\hyperlink{ref-14}{{[}14{]}} & 后缀注意力内核 & 后缀KV长度 & 图兼容依赖特定后缀内核；应用需给定共享关系 \\
动态编译：Nimble（2021，支撑）\protect\hyperlink{ref-15}{{[}15{]}} & 字节码与候选内核 & 形状、分支、内核选择 & 形状函数与VM执行 \\
主机流程编译：DISC（支撑）\protect\hyperlink{ref-35}{{[}35{]}} & 主机与设备代码 & 形状及缓冲区需求 & 静态秩；运行期管理 \\
上界规划：Relax（含CUDA Graph）\protect\hyperlink{ref-10}{{[}10{]}} & 存储及合格子图 & 符号尺寸取值 & 已知上界；预留容量 \\
结构预定：Sequoia（CUDA Graph）\protect\hyperlink{ref-16}{{[}16{]}} & 预选投机树的执行 & token及验证状态 & 结构选择约束 \\
数据复制：xLLM（ACLGraph）\protect\hyperlink{ref-5}{{[}5{]}} & 图内固定偏移布局 & 外部输入内容与地址 & 输入输出复制 \\
计划更新：FlashInfer（CUDA Graph）\protect\hyperlink{ref-13}{{[}13{]}} & 网格与工作区地址 & 工作队列及归约映射 & 计划生成与传输 \\
地址间接访问：GraCE（CUDA Graph）\protect\hyperlink{ref-6}{{[}6{]}} & 指针入口与任务结构 & 输入实际地址 & 指针传输；改写或前置内核 \\
\bottomrule
\end{longtable}
\endgroup

\hypertarget{section-5}{%
\subsection{5 并行执行与资源协调}\label{section-5}}

执行图减少主机重复提交后，设备仍需安排任务、处理依赖和分配资源。提前记录的计算可能受到分支串行、工作量不均和通信等待的影响，逻辑上可并行的任务也可能争用同一资源。本章按系统作出的三类决策展开：首先确定哪些任务可以并行以及需要等待什么，其次决定任务由哪些执行单元承担，最后组织不同资源之间的共置与流水。单卡或多卡、算子或子算子粒度、编译期或运行期决策，作为比较这些方法的属性，而不另建分类层次。

\hypertarget{section-5-1}{%
\subsubsection{5.1 并行计划与依赖约束}\label{section-5-1}}

并行计划首先需要保持任务之间的先后关系。例如，两个前驱可以同时执行，消费二者输出的后继却必须等待；如果前驱和后继位于不同设备，数据传输也成为依赖的一部分。计划既要保留有效并行机会，又要减少为实现这些机会引入的同步、资源争用和控制开销。

Nimble（Kwon等，2020）从依赖结构出发，将流分配转化为图上的匹配问题，以保持逻辑并发并减少跨流同步。同一流内的顺序满足一部分依赖，其余跨流关系通过事件落实；随后，预执行将流映射和同步一起捕获为CUDA Graph。这样，系统无需逐轮在主机上重排任务，独立分支也能够按既定计划执行。该方法以重复执行相同计算集合的静态网络为对象，解决的是计划复用和依赖表达；实际重叠程度还取决于各任务对设备资源的占用\protect\hyperlink{ref-11}{{[}11{]}}。

IOS将目标设备上的测量引入计划选择。它把计算划分为依次执行的阶段，在阶段内部比较两种组织方式：合并符合条件的同类算子，或将算子分成组，使组内串行、组间并行。阶段延迟由硬件实测获得，再通过动态规划组合整体调度。相较于仅分析依赖结构，IOS能够比较算子合并后的计算效率与并发执行的实际代价。但阶段仍是调度单元，其边界不等于CUDA Graph捕获边界；算子合并也受到类型和实现条件限制\protect\hyperlink{ref-22}{{[}22{]}}。

Opara进一步区分流分配与算子发射次序。流分配尽量将节点放入某个前驱所在的流，以减少同步；发射次序则利用共享内存、寄存器和线程需求，在就绪集合中优先考虑资源需求较小的任务，并交替考虑计算密集与访存密集算子。大任务若提前占满资源，可能阻碍后续任务，即使二者在依赖图中彼此独立。Opara将所选多流计划及必要事件捕获为CUDA Graph，把资源分析与低开销回放连接起来。其启发式搜索给出可实施的计划，但不构成任意图上的全局最优保证\protect\hyperlink{ref-19}{{[}19{]}}。

计划还受到前端计算图改写的影响。融合或替换算子可以缩短局部计算，也可能形成较大的关键路径节点，损失原有并行机会。AutoGraph采用兼顾关键路径的代价估计搜索图替换，再以基于Nimble的多流映射和CUDA Graph执行候选图，用端到端测量筛选后续搜索方向。它所划分的是改写搜索空间，而不是独立捕获区域。结构分析确定哪些顺序合法，设备测量与资源分析判断这些顺序是否高效，图改写则进一步改变需要被调度的任务集合\protect\hyperlink{ref-20}{{[}20{]}}。

Mirage提供了相邻的跨层优化视角。它用kernel、thread block和thread三个层次的图表示张量程序，联合探索代数变换、调度变换及新内核生成；候选通过等价性检查后，再优化布局、执行顺序与中间存储。在生成内核内部，算子顺序会影响线程同步位置；在内核之间，融合和存储安排会改变设备内存访问。其评测使用CUDA Graph降低各基线的启动开销，说明研究比较的是这一提交条件下仍然存在的计算组织差异。Mirage的多层图是编译表示，适合用于解释计划与内核实现的联系，而不是图捕获机制本身\protect\hyperlink{ref-40}{{[}40{]}}。

当任务分布到多个设备，计划还必须让远端依赖及时生效。Mustard的静态场景预先确定任务所属GPU和数据传输目标，并向图中加入依赖更新、依赖等待及数据传输节点。本地依赖继续由CUDA Graph管理，跨节点依赖通过NVSHMEM原子操作更新共享状态；接收方在所需前驱完成后继续执行。节点内与节点间数据传输使用对应的通信路径。这里任务归属已经固定，改进集中于依赖推进，减少跨设备协调对主机控制的依赖。它与运行时决定就绪任务归属解决的是不同决策，后者将在下一节讨论\protect\hyperlink{ref-21}{{[}21{]}}。

这些方法共同说明，并行计划不只是为节点排列一个顺序。它还涉及任务能否合并、使用哪些执行流、同步应插在哪里，以及远端数据何时到达。静态计划能复用这些决策，但计划的有效性和收益仍依赖被保留的结构、资源与数据放置条件。

\hypertarget{section-5-2}{%
\subsubsection{5.2 任务分配与负载均衡}\label{section-5-2}}

在依赖关系允许执行之后，系统还要决定由哪个执行单元承担任务。当算子工作量差异较大，或执行时间随输入变化时，仅暴露并行关系可能留下长尾任务和空闲资源。任务分配可以提前完成，也可以在运行中调整；细分任务增加调整空间，同时增加任务描述、队列和依赖管理成本。

Rammer以rTask表示可独立执行的细粒度计算，将算子组织为一组同质任务。例如，矩阵乘法可按输出分块形成多个rTask，不同分块实现又产生资源使用和执行时间不同的候选内核。编译器利用这些信息，为虚拟执行单元生成静态rProgram，并以细粒度屏障表达依赖。Rammer同时管理算子内和算子间的分配，避免两层独立决策限制整体计划；若合并后的安排不利，也可以将计算划分为多个rProgram，而非强制整个网络使用一次设备调用\protect\hyperlink{ref-23}{{[}23{]}}。

这种提前分配要求能够获得任务结构及其性能特征。Rammer依赖可用实现或外部工具将算子拆成rTask，并不自动把任意封闭内核转换为可细分任务。对于结构和耗时较稳定的计算，静态分配减少了运行时选择；若耗时显著变化，固定分配则可能出现不均衡。因此，任务粒度与决策时机应分别比较：更小的任务不等于必须动态分配，动态分配也只有在收益覆盖管理成本时才有价值\protect\hyperlink{ref-23}{{[}23{]}}。

MPK把这部分决策放到持久mega-kernel内部。其tGraph用事件连接计算与通信任务，设备worker从队列取任务，scheduler根据事件就绪状态派发后续工作。它结合两种发射方式：提前发射将任务预先放入指定worker队列，等待依赖满足；即时发射在事件就绪后再分配任务，使先空闲的worker接续工作。编译器依据算子耗时是否具有数据依赖性标记相应区域，并在全局屏障消除累积不均衡后恢复提前发射。这样，稳定区域可以减少调度，不稳定区域则保留在就绪后决定分派位置的机会\protect\hyperlink{ref-24}{{[}24{]}}。

Rammer与MPK都减少逐算子主机决策，但前者强调编译期分配，后者保留事件驱动的设备端调度。MPK的worker、scheduler和队列本身也占用资源，即时发射增加通信与同步，提前发射则降低负载适应能力。二者是理解任务分配的相关执行模型：rProgram和持久内核中的tGraph并不等同于CUDA Graph实例，不能把相应收益直接归入图回放接口的改进\protect\hyperlink{ref-23}{{[}23}, \protect\hyperlink{ref-24}{24{]}}。

Mustard的节点内动态场景则使用CUDA Graph子图作为分配单位。系统将静态任务图拆成可由设备启动的子图，提前在各GPU上初始化；常驻设备侧调度器从共享就绪队列取任务，结合资源占用信息启动子图。插入图内的依赖更新节点在具体任务完成时释放后继工作，无需总是等到所属整个子图结束。变化的是任务执行位置和启动时机，原有依赖拓扑保持固定；子图内部仍由CUDA Graph管理，子图之间再由设备端协调\protect\hyperlink{ref-21}{{[}21{]}}。

这一方案受到数据与设备侧图启动能力的共同限制。Mustard将跨节点执行限定于不采用动态负载均衡的静态分配；动态场景面向数据能容纳于单GPU内存的应用，尚未处理一般的数据所有权迁移与拓扑感知放置。设备侧启动数量限制可能迫使系统合并任务，占用信息也不足以完整描述寄存器和缓存需求。因此，同属多GPU执行，静态依赖推进与动态负载均衡仍有不同的适用边界。MPK以更细的计算和通信任务响应数据到达事件，Mustard则保留现有CUDA Graph子图，两者需要在任务粒度、可复用内核和调度成本上比较，而不能仅比较支持的设备数量\protect\hyperlink{ref-21}{{[}21}, \protect\hyperlink{ref-24}{24{]}}。

\hypertarget{section-5-3}{%
\subsubsection{5.3 资源共置与流水组织}\label{section-5-3}}

任务已经分配之后，还需要决定哪些工作同时占用设备、各自使用多少资源，以及不同阶段如何交错推进。计算、访存与通信可能具有互补的瓶颈，但并发任务仍会共同使用线程、寄存器和缓存。资源互补需要通过任务拆分、共置机制及干扰控制落实；增加并行度也可能破坏原有批量计算效率。

NanoFlow以nano-batch拆分操作，使不同子批次在依赖许可时交错推进。搜索先依据无干扰的内核性能确定nano-operation的数量、批量和顺序，形成减少计算空隙的流水；再固定这些结构性选择，利用实际内核干扰测量调整资源分配和实现选择。这样的分解降低了搜索复杂度，也保留了并发时重新衡量内核组合的必要性。系统还考察通信操作的等价变换，把通信形式对流水依赖的影响纳入计划\protect\hyperlink{ref-25}{{[}25{]}}。

NanoFlow也将设备执行与主机准备重叠：当前迭代尚未结束时，运行时已开始形成后续批次。这可以隐藏CPU工作，却意味着不能立即根据当前步的结束标记移除请求。该系统假设实例获得充足请求，以维持有效批量；当模型或输入、输出长度等工作负载显著变化时，需要重新搜索流水。其方法体现的是持续负载下结构选择、资源分配与状态可见性之间的权衡，不能将吞吐收益视为任意到达强度下的固定性质\protect\hyperlink{ref-25}{{[}25{]}}。

POD-Attention将问题收敛到混合批次中的prefill与decode attention。仅把两类计算提交到不同流，无法保证它们在同一流式多处理器上共置；过细的线程束级融合又可能让较快的一部分等待另一部分。POD-Attention采用线程块级融合，让线程块实际到达流式多处理器后读取其标识SMID，并结合该SM的计数器决定执行prefill还是decode。不同线程块能够独立完成，同步屏障的影响限制在线程块内部；软件调度由此控制共置比例，同时减少一类计算对另一类计算的牵制\protect\hyperlink{ref-26}{{[}26{]}}。

昇腾上的xLLM在算子层讨论Cube与Vector单元的资源分配。它根据负载为并发算子分配不同数量的计算单元，以缩小完成时间差，减少一种资源工作时另一种资源空闲的情况。与NanoFlow以实测干扰修正流水不同，xLLM的资源分配公式采用简化模型，假设算子同时开始，并忽略数据依赖和通信延迟。该机制是昇腾异构资源协调的直接案例，但目标函数不等同于完整解码迭代的最短执行时间，其贡献也需与ACLGraph捕获和回放分别辨识\protect\hyperlink{ref-5}{{[}5{]}}。

三项工作的共同决策是如何让互补资源在时间上重叠，控制层次则不同：NanoFlow搜索跨操作、跨子批次流水，POD-Attention改造单个注意力内核，xLLM协调异构计算单元。相应前提包括持续批量和干扰测量、专用内核实现，以及简化资源模型。表4将这些方法与前两节的计划和分配方法放在共同决策下比较。一次推理可以组合图回放、任务分配和资源重叠，但共享资源的机制需要联合评估，其收益不能直接相加。

\begingroup\small
\begin{longtable}[]{@{}>{\raggedright\arraybackslash}p{(\linewidth - 6\tabcolsep) * \real{0.2}}>{\raggedright\arraybackslash}p{(\linewidth - 6\tabcolsep) * \real{0.22}}>{\raggedright\arraybackslash}p{(\linewidth - 6\tabcolsep) * \real{0.26}}>{\raggedright\arraybackslash}p{(\linewidth - 6\tabcolsep) * \real{0.32}}@{}}
\caption{表4　并行执行与资源协调的机制比较}\tabularnewline
\toprule
\textbf{工作与机制} & \textbf{主要决策} & \textbf{决策依据} & \textbf{执行范围与条件} \\
\midrule
\endfirsthead
\toprule
\textbf{工作与机制} & \textbf{主要决策} & \textbf{决策依据} & \textbf{执行范围与条件} \\
\midrule
\endhead
Nimble（2020）\protect\hyperlink{ref-11}{{[}11{]}} & 流映射与同步 & 依赖结构、匹配 & CUDA Graph；静态计算集合 \\
IOS \protect\hyperlink{ref-22}{{[}22{]}} & 合并或并行分组 & 阶段延迟实测 & 算子计划；阶段非捕获边界 \\
Opara \protect\hyperlink{ref-19}{{[}19{]}} & 流分配与发射次序 & 资源需求、算子类型 & CUDA Graph；启发式计划 \\
AutoGraph \protect\hyperlink{ref-20}{{[}20{]}} & 改写后的并行计划 & 关键路径、端到端测量 & CUDA Graph执行候选 \\
Mirage \protect\hyperlink{ref-40}{{[}40{]}} & 跨层表示与顺序优化 & 等价检查、布局与存储 & 相关编译；非图捕获机制 \\
Mustard：静态场景 \protect\hyperlink{ref-21}{{[}21{]}} & 跨设备依赖推进 & 固定归属、数据传输关系 & CUDA Graph；支持跨节点 \\
Rammer \protect\hyperlink{ref-23}{{[}23{]}} & rTask静态分配 & 编译期性能与资源 & rProgram；需可拆分实现 \\
MPK \protect\hyperlink{ref-24}{{[}24{]}} & 提前或即时派发 & 事件就绪、耗时动态性 & 持久内核；细粒度任务 \\
Mustard：动态场景 \protect\hyperlink{ref-21}{{[}21{]}} & 子图执行位置与时机 & 就绪队列、资源占用 & 节点内；数据可容纳于单GPU \\
NanoFlow \protect\hyperlink{ref-25}{{[}25{]}} & 流水结构与资源配额 & 批量效率、干扰实测 & nano-batch；持续负载条件 \\
POD-Attention \protect\hyperlink{ref-26}{{[}26{]}} & prefill/decode共置 & SM标识与计数器 & 注意力线程块级融合 \\
xLLM资源协调 \protect\hyperlink{ref-5}{{[}5{]}} & Cube/Vector分配 & 完成时间对齐 & 昇腾；模型不计依赖和通信 \\
\bottomrule
\end{longtable}
\endgroup

\hypertarget{section-6}{%
\subsection{6 图实例准备缓存与恢复}\label{section-6}}

前文讨论了图的构建、动态输入适配与并行执行。实际服务还需要回答图何时准备、哪些准备结果值得保留，以及进程变化后怎样继续使用。预热、捕获和实例化需要时间，多个变体也会占用内存；若进程很快缩容或迁移，准备成本可能尚未被回放摊销。因此，本章区分三个决策：当前环境中何时创建和缓存实例，跨进程时如何恢复有效绑定，以及相似实例或环境之间能否消除重复准备。它们与每轮更新活跃图的数据和参数相联系，但管理对象和时间尺度不同。

\hypertarget{section-6-1}{%
\subsubsection{6.1 图实例创建与缓存管理}\label{section-6-1}}

提前捕获让请求能够直接使用已有图，却要求系统预先确定覆盖范围，并为可能很少使用的规格付出准备和存储成本。按需捕获将这部分工作推迟到规格实际出现之后，又会增加首次未命中的处理时间。因而，图缓存不仅承担从已有实例中查找匹配图的职责，还需要决定哪些图提前创建、未命中时如何继续执行，以及保留资源不足时怎样回收。

Hybrid JIT-CUDA Graph将捕获嵌入推理过程。其算法在初始化阶段预捕获短序列图，随后使用不同CUDA流进行捕获和回放。缓存命中时直接回放；未命中时先用JIT路径执行静态计算段，同时异步捕获新图，再将其加入滚动缓存，并在需要时淘汰较少使用的图。原型还让静态图共享激活工作区，并通过事件协调执行路径，以减少重复分配并保持依赖顺序。该方案把一次未命中与后续复用联系起来，使请求不必全部等待新图准备完成。这里依据的是arXiv:2604.23467v1所描述的实现\protect\hyperlink{ref-27}{{[}27{]}}。

后台准备并不消除准备成本。该文指出，所用PyTorch与cuBLAS集成仍会串行化部分捕获操作，每个新序列长度也需要首次捕获。缓存限制保留数量后，再次出现的旧规格还可能重新触发准备。其评测采用单GPU、批量为1的LLaMA-2 7B，并使用预热缓存，因而不能仅据命中后的延迟推断空缓存启动、变化批量或多GPU负载下的收益。比较这类策略，需要同时记录首次未命中、后台干扰、重复命中和资源占用\protect\hyperlink{ref-27}{{[}27{]}}。

在昇腾上，xLLM的自适应图模式结合参数化与多缓存部分图，以有限数量的预编译结果服务后续请求。该设计提供了ACLGraph多图复用的直接案例，但报告没有说明在线淘汰规则或后台捕获算法。xLLM说明需要准备什么实例，Hybrid进一步讨论实例怎样随运行过程产生和保留；二者可按形状覆盖、准备次数和缓存占用比较，不能据前者的多图支持补写相同的在线管理器。当前进程内缓存解决的是实例仍然存在时的复用，跨进程保留准备结果还需要处理另一组绑定问题\protect\hyperlink{ref-5}{{[}5}, \protect\hyperlink{ref-27}{27{]}}。

\hypertarget{section-6-2}{%
\subsubsection{6.2 跨进程图恢复与地址策略}\label{section-6-2}}

图节点中的设备指针和内核函数引用，通常关联原进程的内存分配与模块加载结果。新进程得到相同拓扑，并不意味着这些引用仍然有效。跨进程恢复因此至少需要重新建立任务结构、数据地址、内核绑定及必要内容之间的一致关系。Medusa和Foundry都处理这一问题，但对地址采取了不同策略：前者记录指针来源并在恢复时填入新地址，后者重现图所要求的虚拟地址布局，使保存的参数重新有效\protect\hyperlink{ref-28}{{[}28}, \protect\hyperlink{ref-29}{29{]}}。

Medusa在离线阶段记录图节点与参数，以及缓冲区分配、释放和内核启动序列，建立数据指针与分配次序的关联。保存时用分配索引表达指针来源，在线恢复时重放分配序列，再把本次分配的地址填回相应节点。例如，某参数原来指向离线第若干次分配的缓冲区，恢复时需要找到对应的新缓冲区，而不要求地址数值与过去相同。释放和再次分配可能复用同一地址，因此该关联还要结合指针使用时的分配历史，不能只按地址数值建立一次映射\protect\hyperlink{ref-28}{{[}28{]}}。

数据地址有效之后，内核函数引用和必要内容仍需恢复。Medusa记录内核符号名称，并在新进程中解析函数；对于延迟加载的模块，通过触发内核促使模块进入当前上下文。针对论文所研究模型，作者利用首层计算触发后续层所需模块的加载，减少重新运行完整预热的需要。缓冲区内容则区别处理：模型加载已恢复的权重与可重新计算的临时内容无需重复保存，而回放所需保留的数据仍需恢复。因此，它保存的不只是拓扑，也包括绑定索引、内核信息和必要数据\protect\hyperlink{ref-28}{{[}28{]}}。

这一路线依赖于能够识别并重新解析需要修补的引用。Medusa原型仅拦截主机侧分配并处理直接指针，设备侧分配、间接指针以及多GPU支持在所读论文中仍有待扩展。内核触发策略也依赖具体模型与库的加载行为。对其恢复能力的判断，应围绕这些已经记录的对象和支持条件展开，而不应将图专项物化等同于整个进程的完整快照\protect\hyperlink{ref-28}{{[}28{]}}。

Foundry采用地址重现来减少逐项修补数据指针的需要。一些库将指针嵌入不透明参数结构，若改变数据位置，识别全部内嵌引用就需要了解这些结构。Foundry拦截虚拟内存管理请求，把分配安排在固定基址后的连续虚拟地址区域；恢复时重新建立相应布局，使图参数中的地址继续指向正确位置。仅固定基址还不够：跳过捕获后，捕获窗口中原本发生的临时分配不会自然出现，因此系统也记录并重放这段分配序列，并依据离线地址范围预先映射所需区域。后续分配由此能在已经准备的空间内推进\protect\hyperlink{ref-29}{{[}29{]}}。

地址重现同样不能代替内核恢复。Foundry记录实际使用的模块或库二进制、加载选项和函数目录，恢复时重新加载并解析函数句柄；需要设备侧初始化的模块还要执行相应准备。与Medusa依靠已记录的符号及触发策略建立内核引用相比，Foundry将所需二进制纳入物化对象。两种策略的差异落在保存什么、恢复时重新解析什么，以及必须保持哪些条件；不能据此把二者排列成对所有模型都成立的优劣顺序\protect\hyperlink{ref-28}{{[}28}, \protect\hyperlink{ref-29}{29{]}}。

图3概括了两条地址路线。重绑定允许恢复后的数据地址改变，但要求找全并修补需要更新的引用；地址重现减少数据指针修补，却要求可重建一致的虚拟布局及分配过程。二者都需要处理内核与必要数据，并且都要在恢复后检验依赖和参数是否仍然有效。

\begin{figure}
\centering
\includegraphics[width=\linewidth]{figures/restoration.png}
\caption{\textbf{图3　跨进程恢复的地址策略。左侧根据分配索引重绑定到新地址，右侧重现预期虚拟布局；图为对Medusa与Foundry机制的归纳示意，不表示全部恢复步骤或耗时比例}\protect\hyperlink{ref-28}{{[}28}, \protect\hyperlink{ref-29}{29{]}}\textbf{。}}
\end{figure}

通用GPU恢复研究提供了更大状态范围的对照。CRIUgpu通过驱动插件协调CPU与GPU状态：先锁定相关执行并保存状态，恢复原地址映射、流和上下文后再恢复运行；其论文所用版本尚不支持NCCL checkpoint。GCR则将控制状态交由驱动恢复，单独处理数据buffer：保存原虚拟地址关系，释放物理内存，恢复时重新映射数据。GCR原型未实现并发C/R。这些方法说明地址一致性也可以由进程或驱动层维持，但它们的保存对象、停顿条件和恢复代价与图专项物化不同；上述原文不应被扩展为独立图模板可序列化的证明\protect\hyperlink{ref-41}{{[}41}, \protect\hyperlink{ref-42}{42{]}}。

\hypertarget{section-6-3}{%
\subsubsection{6.3 图模板共享与重复准备消除}\label{section-6-3}}

单张图能够恢复后，大量相似实例仍可能重复构建相同结构、加载相同环境。消除这类工作首先要明确共享单位：多个规格可以共享图模板，多个rank可以在满足条件时共享离线准备结果，服务实例则可以复用预热环境或模型数据。共享对象不同，需要维持的一致性也不同。

Foundry依据节点类型、依赖及相关属性构造拓扑标识，将结构相同的图归为一组。每组建立一个模板，其他规格保存各自参数集；执行目标规格前，通过图更新机制应用相应参数，避免为每种规格重新构建完整实例。这并不要求所有输入使用同一批量：能够共享的是符合结构兼容条件的模板，批量对应的参数仍需更新。模板复用因而把重复的结构准备与必要的规格变化分开，在线成本则从完整构建转为相应参数更新\protect\hyperlink{ref-29}{{[}29{]}}。

在多rank场景中，Foundry进一步利用一致执行流共享准备结果。离线阶段使用通信占位实现捕获，在线替换真实通信函数并补入rank相关参数。这依赖各rank具有等价图结构，不覆盖执行不同模型阶段的流水线并行。其vLLM集成还固定KV缓存大小，以保持不同运行之间的布局一致。因此，多规格和多rank共享都需要明确兼容条件：结构相同不意味着地址、内核或通信状态天然相同，减少模板数量也不意味着任意配置可以使用同一存档。以上分析限于arXiv:2604.06664v1所描述和验证的范围\protect\hyperlink{ref-29}{{[}29{]}}。

图之外的初始化也能形成重复准备。DEEPSERVE在昇腾上保留预热Pod与执行引擎，将Python库、NPU状态和HCCL连接初始化移出扩容关键路径；NPU-fork再从已有引擎传输模型权重。这两项机制分别复用环境与模型数据，能够支撑服务实例更快可用，但论文没有据此证明已捕获ACLGraph实例可以跨进程恢复。将其作为支撑方法，能够补齐图恢复之后仍然需要满足的服务条件，而不会混淆恢复对象\protect\hyperlink{ref-30}{{[}30{]}}。

PhoenixOS也复用部分环境准备。常驻进程预先建立CUDA、cuBLAS上下文及部分通信资源，恢复时分配已有上下文，后续访问需要经进程映射调用。它还明确支持图内核的读写追踪，用于其GPU检查点恢复机制；这一能力与保存图节点并恢复跨进程绑定仍需分别说明。环境池可以减少在线初始化，却带来常驻资源和访问路径的代价，不能仅以启动等待减少推断完整运行成本同等下降\protect\hyperlink{ref-31}{{[}31{]}}。

表5按复用对象和在线剩余工作比较上述方法。图缓存、地址恢复、模板共享和环境预热可以作用于不同阶段；组合评价应分别记录离线准备、常驻资源、首次可用时间及后续调用成本，再观察完整服务收益。对于ACLGraph，xLLM提供了多图复用的直接材料，DEEPSERVE说明环境准备具有独立价值；跨进程重绑定、地址重现与模板共享仍需根据目标版本的接口和运行配置逐项验证，不能由GPU上的机制比较直接推出平台支持结论。

\begingroup\small
\begin{longtable}[]{@{}>{\raggedright\arraybackslash}p{(\linewidth - 8\tabcolsep) * \real{0.18}}>{\raggedright\arraybackslash}p{(\linewidth - 8\tabcolsep) * \real{0.17}}>{\raggedright\arraybackslash}p{(\linewidth - 8\tabcolsep) * \real{0.23}}>{\raggedright\arraybackslash}p{(\linewidth - 8\tabcolsep) * \real{0.22}}>{\raggedright\arraybackslash}p{(\linewidth - 8\tabcolsep) * \real{0.2}}@{}}
\caption{表5　图准备与恢复方法的机制比较}\tabularnewline
\toprule
\textbf{工作与机制} & \textbf{复用对象} & \textbf{保存或共享方式} & \textbf{在线剩余工作} & \textbf{主要条件} \\
\midrule
\endfirsthead
\toprule
\textbf{工作与机制} & \textbf{复用对象} & \textbf{保存或共享方式} & \textbf{在线剩余工作} & \textbf{主要条件} \\
\midrule
\endhead
Hybrid JIT-CUDA Graph \protect\hyperlink{ref-27}{{[}27{]}} & 进程内图实例 & 按规格滚动缓存 & 未命中执行与捕获 & 单GPU原型；有捕获干扰 \\
xLLM \protect\hyperlink{ref-5}{{[}5{]}} & 多缓存部分图 & 参数化与有限预编译 & 模式选择、数据准备 & ACLGraph直接；未给淘汰策略 \\
Medusa \protect\hyperlink{ref-28}{{[}28{]}} & 图及必要绑定 & 分配索引、符号与数据 & 新地址重绑定、内核解析 & 主机分配、直接指针、单GPU \\
Foundry：环境恢复 \protect\hyperlink{ref-29}{{[}29{]}} & 图执行环境 & 固定布局、分配序列、二进制 & 地址重现、模块加载 & 布局与配置兼容 \\
Foundry：模板共享 \protect\hyperlink{ref-29}{{[}29{]}} & 同结构图模板 & 共享拓扑、规格参数集 & 参数更新、通信替换 & 跨rank要求等价执行流 \\
CRIUgpu \protect\hyperlink{ref-41}{{[}41{]}} & CPU/GPU进程状态 & 驱动级统一快照 & 映射与控制状态恢复 & 所读版本不支持NCCL C/R \\
GCR \protect\hyperlink{ref-42}{{[}42{]}} & 控制状态与数据 & 驱动与buffer分别处理 & 原虚拟地址下重新映射 & 通用恢复；原型不并发C/R \\
DEEPSERVE \protect\hyperlink{ref-30}{{[}30{]}} & 预热环境、权重 & 环境池、NPU-fork & 组合资源与后续准备 & 昇腾服务支撑 \\
PhoenixOS \protect\hyperlink{ref-31}{{[}31{]}} & 预创建上下文 & 常驻池与状态恢复 & 恢复状态、映射调用 & 通用GPU恢复支撑 \\
\bottomrule
\end{longtable}
\endgroup

\hypertarget{section-7}{%
\subsection{7 综合比较与昇腾适用性}\label{section-7}}

\hypertarget{section-7-1}{%
\subsubsection{7.1 跨类别特征与方法选择}\label{section-7-1}}

四类方法分别改变构建、适配、执行和实例管理决策，收益却会通过共享状态和资源相互影响。扩大捕获范围可以减少区域间的主机提交，也会把更多动态操作放进同一复用约束。xLLM的部分图保留图外路径，GraCE通过转换代码消除部分捕获障碍；二者既可能替代，也可以先后使用。若瓶颈来自区域切换，扩大捕获可能有价值；若瓶颈来自每轮数据搬运，则应先比较状态更新方案。捕获覆盖率本身不能决定这些选择\protect\hyperlink{ref-5}{{[}5}, \protect\hyperlink{ref-6}{6{]}}。

形状专门化与缓存管理具有另一种联系。为不同规格准备图可以减少填充，却提高准备次数和保留资源；FlashInfer的配置选图回答如何匹配已有实例，Hybrid的在线捕获回答未命中时怎样形成后续可复用的实例。选择更细的规格，是否有足够重复次数摊销新增准备，应通过实际输入分布判断。缓存命中率还应与命中后的填充及更新代价一起报告\protect\hyperlink{ref-13}{{[}13}, \protect\hyperlink{ref-27}{27{]}}。

表6以统一字段列出代表机制的横向属性。比较同一列能够看出方法是否依赖相同条件；比较相邻列则能解释哪些机制可以组合。例如FlashInfer保持工作区布局，GraCE保持指针入口，二者都需要正确的更新顺序，但改变的是不同内容。Rammer与MPK都暴露细粒度任务，仍因分配时机不同而具有不同的调度成本。单卡或多卡仅说明执行范围，不能替代对任务所有权与依赖推进方式的分析。

\begingroup\small
\begin{longtable}[]{@{}>{\raggedright\arraybackslash}p{(\linewidth - 8\tabcolsep) * \real{0.18}}>{\raggedright\arraybackslash}p{(\linewidth - 8\tabcolsep) * \real{0.16}}>{\raggedright\arraybackslash}p{(\linewidth - 8\tabcolsep) * \real{0.25}}>{\raggedright\arraybackslash}p{(\linewidth - 8\tabcolsep) * \real{0.21}}>{\raggedright\arraybackslash}p{(\linewidth - 8\tabcolsep) * \real{0.2}}@{}}
\caption{表6　代表机制的横向属性比较}\tabularnewline
\toprule
\textbf{具体机制} & \textbf{主要决策时机} & \textbf{保持稳定与允许变化的对象} & \textbf{控制粒度或执行范围} & \textbf{证据关系} \\
\midrule
\endfirsthead
\toprule
\textbf{具体机制} & \textbf{主要决策时机} & \textbf{保持稳定与允许变化的对象} & \textbf{控制粒度或执行范围} & \textbf{证据关系} \\
\midrule
\endhead
xLLM部分图与固定池\protect\hyperlink{ref-5}{{[}5{]}} & 初始化及每轮模式选择 & 图内布局稳定；输入值和部分维度可变 & 昇腾模型模块 & ACLGraph直接 \\
GraCE候选选择与间接访问\protect\hyperlink{ref-6}{{[}6{]}} & 编译慢路径选优；每轮更新 & 指针入口稳定；外部数据地址可变 & 编译区域与内核参数 & CUDA Graph直接 \\
FlashInfer计划与执行分离\protect\hyperlink{ref-13}{{[}13{]}} & CPU逐轮规划；GPU消费 & 工作区地址与启动配置稳定；计划内容可变 & Attention任务队列 & CUDA Graph直接适配 \\
Hydragen后缀内核\protect\hyperlink{ref-14}{{[}14{]}} & 解码迭代 & 实现保持图兼容；后缀KV长度可变 & Attention子计算 & 图兼容使用实例 \\
Rammer静态rProgram\protect\hyperlink{ref-23}{{[}23{]}} & 编译期 & 任务计划固定；按输入计算数值 & 细粒度rTask & 相邻执行模型 \\
MPK混合任务发射\protect\hyperlink{ref-24}{{[}24{]}} & 编译期划区及运行期分派 & 事件依赖已描述；任务分配可动态调整 & 持久内核内计算与通信任务 & 相邻执行模型 \\
NanoFlow子批次流水\protect\hyperlink{ref-25}{{[}25{]}} & 离线搜索及在线批量形成 & 选定流水与资源配置；请求内容变化 & 子批次及计算通信操作 & 系统调度支撑 \\
Mustard跨节点静态场景\protect\hyperlink{ref-21}{{[}21{]}} & 预先分配；设备端推进依赖 & 任务归属固定；就绪状态变化 & 多GPU任务图 & CUDA Graph直接 \\
Mustard节点内动态场景\protect\hyperlink{ref-21}{{[}21{]}} & 设备运行期 & 预建子图；执行位置按就绪任务分配 & 单节点多GPU子图 & CUDA Graph直接 \\
Medusa图物化\protect\hyperlink{ref-28}{{[}28{]}} & 离线记录及在线恢复 & 保存分配关联；新进程实际地址可变化 & 图节点及绑定 & CUDA Graph直接 \\
Foundry环境与模板物化\protect\hyperlink{ref-29}{{[}29{]}} & 离线物化及在线重建 & 重现地址布局；同拓扑规格参数可变 & 环境、模板及多rank参数 & CUDA Graph直接 \\
\bottomrule
\end{longtable}
\endgroup

方法组合首先应检查共享状态是否兼容。固定工作区节约分配，但多个回放若共享同一位置，就必须保证前一轮读取完成之后再写入下一轮数据。固定地址不是访问顺序的保证，资源重叠也不自动遵守缓存对象的生命期。与此同时，图恢复加快只能减少它所覆盖的启动阶段；当权重或通信初始化仍占主导时，需要与DEEPSERVE等环境准备机制联合分析，不能将图阶段加速比直接当作首次可服务时间的加速比\protect\hyperlink{ref-5}{{[}5}, \protect\hyperlink{ref-13}{13}, \protect\hyperlink{ref-28}{28}, \protect\hyperlink{ref-29}{29}, \protect\hyperlink{ref-30}{30{]}}。

\hypertarget{section-7-2}{%
\subsubsection{7.2 昇腾工程证据与迁移条件}\label{section-7-2}}

xLLM给出了ACLGraph部分捕获、动态维度和固定偏移内存池的直接论文证据；其图模式消融可与整体系统评测区分讨论。Cube与Vector资源协调属于同一系统中的另一机制，不能用捕获收益证明资源分配收益，反之亦然。DEEPSERVE说明环境预热和引擎准备能够改善昇腾服务启动，但其结论也不能替代已捕获图实例的跨进程恢复证据\protect\hyperlink{ref-5}{{[}5}, \protect\hyperlink{ref-30}{30{]}}。

公开工程材料补充了论文没有展开的接口分工。Ascend/pytorch的已合并PR !20264新增任务组捕获、任务组参数更新接口，以及用于更新时序控制的ExternalEvent。它直接说明运行时具有组织指定任务参数更新的工程入口，未给出任意算子均可更新或任意并发组合均正确的证明\protect\hyperlink{ref-43}{{[}43{]}}。vLLM-Ascend v0.13.0文档进一步把分桶填充、超阈值回退、注意力参数更新及流资源约束连接到模型执行。其描述具有版本和配置条件，文档中的模式限制也不能直接推广为所有后续版本的能力边界\protect\hyperlink{ref-44}{{[}44{]}}。

为核实框架与图对象的分工，本文同时检查v0.13.0标签对应提交6281c120中的ACLGraphWrapper实现。该包装器按批次描述符缓存图对象，运行模式不匹配时调用原执行路径；它明确把持久输入缓冲区和输入复制交由外部负责，并在调试模式检查回放地址。其注意力更新路径重新提供序列长度等参数，结束任务更新后记录事件；回放路径通常先执行设备同步，仅对源码限定的完整图草稿模型路径豁免；注释将异步调度或多线程下的更新时序列为同步动机。由此可见，分桶选择、缓冲区管理和更新顺序由多个层次共同完成，不能只依据NPUGraph对象存在就推断状态复用已经正确\protect\hyperlink{ref-45}{{[}45{]}}。

表7将可确认动作与仍需验证的条件分开。表中``有实现''仅表示所引版本中存在该机制；``有限''表示依赖算子、模式或资源条件；``待核实''表示现有证据不足以支持同等能力，不表示平台不存在相应实现。

\begingroup\small
\begin{longtable}[]{@{}>{\raggedright\arraybackslash}p{(\linewidth - 6\tabcolsep) * \real{0.19}}>{\raggedright\arraybackslash}p{(\linewidth - 6\tabcolsep) * \real{0.36}}>{\raggedright\arraybackslash}p{(\linewidth - 6\tabcolsep) * \real{0.15}}>{\raggedright\arraybackslash}p{(\linewidth - 6\tabcolsep) * \real{0.3}}@{}}
\caption{表7　昇腾已有证据及GPU机制的迁移边界}\tabularnewline
\toprule
\textbf{问题与状态} & \textbf{昇腾可确认的证据} & \textbf{GPU对照机制} & \textbf{仍需确认的条件} \\
\midrule
\endfirsthead
\toprule
\textbf{问题与状态} & \textbf{昇腾可确认的证据} & \textbf{GPU对照机制} & \textbf{仍需确认的条件} \\
\midrule
\endhead
区域捕获与参数化：有限 & xLLM部分图及指定算子改造\protect\hyperlink{ref-5}{{[}5{]}} & GraCE兼容性变换、Relax存储规划\protect\hyperlink{ref-6}{{[}6}, \protect\hyperlink{ref-10}{10{]}} & 所用算子、内存规划及图外边界 \\
变体选择与回退：有实现 & 指定版本按批次描述符缓存及按模式分派；文档分桶与阈值回退\protect\hyperlink{ref-44}{{[}44}, \protect\hyperlink{ref-45}{45{]}} & FlashInfer配置选图\protect\hyperlink{ref-13}{{[}13{]}} & 动态输入覆盖、填充和流资源成本 \\
任务参数更新：有限 & 任务组更新接口与ExternalEvent；注意力更新路径\protect\hyperlink{ref-43}{{[}43}, \protect\hyperlink{ref-45}{45{]}} & FlashInfer元数据、GraCE间接访问\protect\hyperlink{ref-6}{{[}6}, \protect\hyperlink{ref-13}{13{]}} & 可更新对象、共享缓冲区生命期和多流顺序 \\
在线实例管理：部分证据 & xLLM多图缓存；包装器保存批次条目\protect\hyperlink{ref-5}{{[}5}, \protect\hyperlink{ref-45}{45{]}} & Hybrid在线捕获与缓存\protect\hyperlink{ref-27}{{[}27{]}} & 未据此建立后台捕获与淘汰算法 \\
跨进程图恢复：待核实 & DEEPSERVE支持环境和权重复用，未证明同等图恢复能力\protect\hyperlink{ref-30}{{[}30{]}} & Medusa重绑定、Foundry布局重现\protect\hyperlink{ref-28}{{[}28}, \protect\hyperlink{ref-29}{29{]}} & 图绑定记录、内核恢复、布局及版本失效条件 \\
资源与多设备协调：有限 & xLLM异构单元协调；图资源预算与通信域相关\protect\hyperlink{ref-5}{{[}5}, \protect\hyperlink{ref-44}{44{]}} & NanoFlow流水、Mustard跨设备依赖\protect\hyperlink{ref-21}{{[}21}, \protect\hyperlink{ref-25}{25{]}} & 资源分配接口、通信语义和具体并行策略 \\
\bottomrule
\end{longtable}
\endgroup

这些工程证据同时影响比较方法。共享输入缓冲区减少内存占用，却可能要求额外同步；较少图桶节约流资源，却可能增加填充；通信域和分段数量变化又会改变可保留图的数量。因此，GPU方法向昇腾迁移时应先明确实现所需能力，再在固定软件版本和负载中测量额外更新、同步及资源占用。公开PR或源码能够说明实现路径，不能充当缺失的性能实验。

\hypertarget{section-7-3}{%
\subsubsection{7.3 原论文评测设置与结果边界}\label{section-7-3}}

执行图研究的测量对象具有不同层次：编译器通常测程序或模型延迟，Attention系统测内核或解码吞吐，任务运行时测DAG完成时间，服务平台则关注启动和请求时延。统一比较的第一步是界定结果所覆盖的工作。首token时延（time to first token，TTFT）包含排队、预填充与可能的准备，每输出token时间（time per output token，TPOT）更接近持续生成成本；图准备时间、图恢复时间与整个进程首次可服务时间也不是同一指标。

表8汇总代表论文已经采用的评测条件，表9记录可定位的结果和归因边界。数值均来自所标版本的原论文，不代表本文重新复现；表中明确每个数字所覆盖的计时范围；未纳入本表的准备成本不能据此视为零，也不据此判定原论文没有报告。不同硬件、输入和基线下的结果用于回答各自研究问题，不按加速比大小排序。对整体系统结果，只有具有相应消融或对照的部分才用于解释图机制贡献。

\begingroup\small
\begin{longtable}[]{@{}>{\raggedright\arraybackslash}p{(\linewidth - 8\tabcolsep) * \real{0.19}}>{\raggedright\arraybackslash}p{(\linewidth - 8\tabcolsep) * \real{0.22}}>{\raggedright\arraybackslash}p{(\linewidth - 8\tabcolsep) * \real{0.22}}>{\raggedright\arraybackslash}p{(\linewidth - 8\tabcolsep) * \real{0.24}}>{\raggedright\arraybackslash}p{(\linewidth - 8\tabcolsep) * \real{0.13}}@{}}
\caption{表8　代表工作的评测设置及证据边界}\tabularnewline
\toprule
\textbf{工作与证据对象} & \textbf{模型与设备} & \textbf{动态性及规模} & \textbf{热冷状态与准备成本} & \textbf{原文位置（PDF页）} \\
\midrule
\endfirsthead
\toprule
\textbf{工作与证据对象} & \textbf{模型与设备} & \textbf{动态性及规模} & \textbf{热冷状态与准备成本} & \textbf{原文位置（PDF页）} \\
\midrule
\endhead
xLLM：ACLGraph直接 \protect\hyperlink{ref-5}{{[}5{]}} & Qwen3-1.7B/4B；Ascend & 表8输入/输出均2048；卡数未单列 & 图模式消融；所列消融未单列准备耗时 & p.24 §5；p.30 §5.2；p.32 表8 \\
GraCE：CUDA Graph直接 \protect\hyperlink{ref-6}{{[}6{]}} & 25项训练/推理/HPC；主H100 NVL & 主实验单卡、指定batch；另TP1/2/4 & 100次运行均值；回放与编译分开；编译含捕获 & pp.11--14 §5、图10--13 \\
Nimble2020：CUDA Graph直接 \protect\hyperlink{ref-11}{{[}11{]}} & CNN；1×V100；推理batch=1 & 给定形状预运行；单GPU多流 & 准备后复用；本表不含准备时间 & pp.8--9 §5、表1 \\
Nimble2021：相邻动态编译 \protect\hyperlink{ref-15}{{[}15{]}} & LSTM/Tree-LSTM/BERT；单T4等 & 动态主实验；表4固定BERT长度128 & 运行时分解；本表不含编译时间 & pp.8--10 §5、表4 \\
FlashInfer：图兼容后端 \protect\hyperlink{ref-13}{{[}13{]}} & Llama3.1-8B/70B；1/4×H100 & ShareGPT/变长请求；约束P99 TTFT\textless200ms & 初始化JIT并缓存；本表不含初始化时间 & p.8 Listing 1、§4.1、图7 \\
Hydragen：图兼容实现，前缀共享评测 \protect\hyperlink{ref-14}{{[}14{]}} & CodeLlama-7B；8×A10040GB & APPS两层/单层共享；每题128候选 & 所选APPS实验排除分词及反分词；本表不含其他准备时间 & pp.10--11 §4.4；p.23 D.4 \\
Rammer：相邻执行调度 \protect\hyperlink{ref-23}{{[}23{]}} & 6类网络；主单V100 & 固定模型；batch=1/4/16 & 1000次均值；本表不含编译及profile时间 & pp.10--12 §5 \\
MPK：相邻设备侧运行 \protect\hyperlink{ref-24}{{[}24{]}} & Qwen3/Llama；A100/H100/B200 & 离线批量；输入64/输出1024；另多卡TP & 运行时延；本表不含预热和编译时间；预印本 & pp.9--12 §6 \\
NanoFlow：相邻资源调度 \protect\hyperlink{ref-25}{{[}25{]}} & 主Llama2-70B；8×A10080GB & 充足离线请求及在线trace；图9固定长度 & 稳态吞吐；搜索在部署准备阶段，耗时不纳入本表 & pp.11--14 §4.2、§6 \\
Mustard：CUDA Graph直接，非LLM \protect\hyperlink{ref-21}{{[}21{]}} & 随机DAG、LU/Cholesky；2/4/8×A100 & 动态任务数据需能放入单卡；跨节点测静态图 & 5次均值；本表不含应用准备时间 & pp.8--11 §5、表2--4 \\
Hybrid：CUDA Graph直接 \protect\hyperlink{ref-27}{{[}27{]}} & Llama2-7B；1×H100；batch=1 & 长度10--500；图缓存 & 暖启动缓存、1000次；不能视为冷捕获；v1预印本 & pp.5--7 §IV--VI、表I--II \\
Medusa：CUDA Graph直接恢复 \protect\hyperlink{ref-28}{{[}28{]}} & 10个0.5B--14B模型；单A100 & 每模型35张批量图；trace为多卡独立实例 & 在线加载不计离线物化；离线均值39.2s另报 & pp.9--12 §7、图8--9 \\
Foundry：CUDA Graph直接恢复 \protect\hyperlink{ref-29}{{[}29{]}} & 稠密/MoE；H200主平台；DP/EP & 每配置512张图；批量1--512 & 暖环境、暖编译缓存、已有归档；排除权重；预印本 & pp.8--11 §6、图7--11 \\
DEEPSERVE：昇腾系统相邻证据 \protect\hyperlink{ref-30}{{[}30{]}} & Llama/CodeLlama/Qwen；TP1/4/8 & 预热TE跨节点扩容；HCCS/RoCE & 图10只计权重就位；已有运行TE；不是完整冷启动 & pp.11--13 §6、图9--11 \\
PhoenixOS：相邻GPU进程恢复 \protect\hyperlink{ref-31}{{[}31{]}} & 训练/推理；单卡或8×A800 & 并发checkpoint/restore；非形状选图 & 主存快照/context池为前提；训练stall与启动另计 & pp.10--13 §8；p.18 表4 \\
\bottomrule
\end{longtable}
\endgroup

表8明确所选结果的计时范围；未纳入本表的准备时间不意味着原论文没有讨论，原文明示排除的成本则不能算入相应结果。尤其是Foundry的服务引擎初始化、Medusa的加载阶段、DEEPSERVE的TE-load与PhoenixOS的checkpoint stall，计时对象并不相同。

\begingroup\small
\begin{longtable}[]{@{}>{\raggedright\arraybackslash}p{(\linewidth - 6\tabcolsep) * \real{0.2}}>{\raggedright\arraybackslash}p{(\linewidth - 6\tabcolsep) * \real{0.29}}>{\raggedright\arraybackslash}p{(\linewidth - 6\tabcolsep) * \real{0.37}}>{\raggedright\arraybackslash}p{(\linewidth - 6\tabcolsep) * \real{0.14}}@{}}
\caption{表9　代表结果、负结果及收益归因}\tabularnewline
\toprule
\textbf{工作与具体条件} & \textbf{原文结果} & \textbf{可支持的结论与边界} & \textbf{原文位置（PDF页）} \\
\midrule
\endfirsthead
\toprule
\textbf{工作与具体条件} & \textbf{原文结果} & \textbf{可支持的结论与边界} & \textbf{原文位置（PDF页）} \\
\midrule
\endhead
xLLM；Qwen3-1.7B；输入/输出各2048 \protect\hyperlink{ref-5}{{[}5{]}} & 关→开自适应图：2385.58→3038.88 token/s；平均TPOT 39.27→30.63ms & 直接图模式消融；不能将收益拆归单项机制；所选表未单列卡数及准备时间 & p.30 §5.2；p.32 表8 \\
GraCE；H100、EOS；无图速度=1 \protect\hyperlink{ref-6}{{[}6{]}} & PyTorch2-CG为0.71×，GraCE为1.00×；全体编译耗时平均为基线2.21× & 图并非总获益；选择性部署避免本例退化，但增加准备成本 & p.12 图10；p.14 图13 \\
Nimble2020；batch=1，单V100 \protect\hyperlink{ref-11}{{[}11{]}} & 多流/自身单流：NASNet-A mobile 1.88×，large 1.31× & 逻辑并发多不必然更快；large的较大核更易占满资源 & p.9 §5.2、表1 \\
Nimble2021；BERT固定长度128，单T4 \protect\hyperlink{ref-15}{{[}15{]}} & TVM 5.58ms；Nimble 5.86ms=核5.60+其他0.26ms & 动态通用性存在静态场景代价；这是编译器对照，非CUDA Graph退化 & p.9 §5.3、表4 \\
FlashInfer；SGLang、8B、1×H100、ShareGPT \protect\hyperlink{ref-13}{{[}13{]}} & 图7所列ITL：Triton 21.7ms→FlashInfer 13.5ms & 后端及动态计划的整体效果；没有隔离CUDA Graph贡献 & p.8 §4.1、图7 \\
Hydragen；CodeLlama-7B、8×A100，APPS \protect\hyperlink{ref-14}{{[}14{]}} & bs=2道题时，两层共享比单层降18\%；再增至8道题额外降45\% & 分开解释计算效率和节省显存后的增批收益；不能把两百分数相加 & pp.10--11 §4.4；p.23 D.4 \\
MPK；Qwen3-8B、单A100、低批量解码 \protect\hyperlink{ref-24}{{[}24{]}} & 逐token时延由vLLM/SGLang的14.5ms降至12.5ms & 整体设备侧运行时收益；不是仅移除launch的消融；离线负载 & p.10 §6 \\
NanoFlow；70B、8×A100；输入512/输出1024 \protect\hyperlink{ref-25}{{[}25{]}} & 每GPU吞吐：不重叠1048，仅拆小批952，完整方案1277 token/s & 只拆分会退化；调度重叠与资源配置协同后才获益 & p.14 §6.4、图9 \\
Mustard；Cholesky 60K，8×A100 \protect\hyperlink{ref-21}{{[}21{]}} & 对单卡库：Mustard 2.728×，细分块cuSmGF 3.185× & 并非所有规模都优于专门库；双方分块不同，且非LLM证据 & p.11 表4、Scaling段 \\
Hybrid v1；H100、batch=1，生成长度500 \protect\hyperlink{ref-27}{{[}27{]}} & P99逐token：Hybrid 23.68ms，HF 18.38ms，TensorRT-LLM 16.65ms & 所读版本表II在长生成区间未显示一致尾延迟优势；不能由短序列推广 & p.7 表II；pp.5--6 设置 \\
Medusa；Qwen1.5-4B、单A100 \protect\hyperlink{ref-28}{{[}28{]}} & 加载2.85→1.67s；图阶段0.90s→预热0.31s+恢复0.26s；KV 0.50→0.02s & 降幅来自多机制及重叠；10个模型离线准备均值39.2s另计；不能全归图恢复 & pp.10--11 §7.3--7.4、图8--9 \\
Foundry；Qwen3-14B、H200、DP1--8 \protect\hyperlink{ref-29}{{[}29{]}} & 引擎初始化：带图vLLM 36--48s→Foundry 1.7--1.8s & 暖环境/编译缓存、已有归档且排除权重；不是完整机器冷启动 & p.9 §6.2.1、图7 \\
\bottomrule
\end{longtable}
\endgroup

这些结果共同支持三点更审慎的比较。首先，回放、并行与动态通用性都有适用条件：GraCE、Nimble2021、NanoFlow和Mustard均提供了不同形式的负例。其次，需要优先使用能隔离机制的对照，例如xLLM图模式开关、Nimble2020自身单流基线，以及Hydragen保持batch不变的消融。最后，准备与恢复不是同一计时过程；Medusa把离线准备另列，Foundry明确排除环境和权重成本，因此这些工作说明``哪些启动阶段可复用''，并不能直接给出统一冷启动排名。

实际证据表可以区分能力覆盖与性能覆盖。一个实现能够处理长度变化，并不说明它已在连续到达、冷缓存和混合预填充/解码下评测；多GPU任务图的依赖推进，也不能直接代表包含KV所有权和请求迁移的LLM服务。对于输入内容相关的路径变化、图缓存失效和跨版本恢复，应分别检查原文是否给出实现、正确性验证及性能测试，不能把其中一项自动替代另外两项。

部署负载还应覆盖改变执行条件的真实来源。S-LoRA同时管理适配器与KV存储，Punica通过分段计算合并多个LoRA模型的执行；两者提示适配器组成、rank与加载状态可能影响可复用配置。这里的rank指低秩维度，与分布式进程rank含义不同。Sarathi-Serve的分块预填充则使每轮的预填充与解码比例成为重要变量。纳入这些负载有助于检验图变体覆盖和状态准备，但它们的系统收益不能归为执行图接口优化\protect\hyperlink{ref-32}{{[}32}, \protect\hyperlink{ref-33}{33}, \protect\hyperlink{ref-34}{34{]}}。

图模式也不能消除所有主机开销。CPU-Induced Slowdowns关注多GPU推理中主机处理和同步的影响；图前后的批量形成、输入准备和跨rank协调仍应单独剖析。The Silent Hyperparameter指出推理后端与配置可能影响输出行为，相关差异不能全部归因于图模式。性能结果因此需要同时交代后端版本、精度与采样条件，并在相同模型设置下检查输出一致性\protect\hyperlink{ref-46}{{[}46}, \protect\hyperlink{ref-47}{47{]}}。

\hypertarget{section-8}{%
\subsection{8 研究挑战与发展方向}\label{section-8}}

\hypertarget{section-8-1}{%
\subsubsection{8.1 面向负载分布的构建与缓存联合决策}\label{section-8-1}}

GraCE通过局部候选实测选择图模式，FlashInfer通过配置与逐轮计划分开处理动态性，Hybrid尝试在未命中后补充捕获结果。这些方法分别优化选择、适配和创建，但其证据窗口不同：编译期选优未必反映后续形状分布漂移，预热后的缓存测试也不能说明首次未命中的尾时延。图实例数量增加还会竞争显存及运行时资源。因此，需要研究在给定资源预算下，如何联合选择捕获区域、规格覆盖、创建时机和保留时间\protect\hyperlink{ref-6}{{[}6}, \protect\hyperlink{ref-13}{13}, \protect\hyperlink{ref-27}{27{]}}。

这一问题应通过同一请求轨迹中的不同策略比较来判断：固定预捕获、按需捕获、负载感知创建和淘汰分别保留多少图，经历多少首次未命中，产生多少填充和后台干扰。平均吞吐之外，应观察P99时延、准备成本摊销轮次以及负载改变后资源回收的速度。若减少图数量同时显著增加填充，或后台捕获使在线尾时延恶化，即使缓存命中率提高，也不足以证明策略有效。

\hypertarget{section-8-2}{%
\subsubsection{8.2 状态更新与并发执行的共同约束}\label{section-8-2}}

xLLM稳定图内位置，FlashInfer稳定工作区布局，GraCE稳定指针入口；三者允许变化的状态不同，却都需要更新与消费的正确顺序。昇腾任务组更新与ExternalEvent提供了表达部分顺序的工程能力，框架回放路径仍可能需要同步协调。单独证明参数可以更新，不能覆盖多轮、多流共享缓冲区时的重叠执行\protect\hyperlink{ref-5}{{[}5}, \protect\hyperlink{ref-6}{6}, \protect\hyperlink{ref-13}{13}, \protect\hyperlink{ref-43}{43}, \protect\hyperlink{ref-45}{45{]}}。

可将每项状态的写入者、读取者、有效期和回收条件显式关联到图计划，并比较共享缓冲、双缓冲或独立工作区的取舍。验证应包含更新早于前轮读取结束、图未命中后工作区复用、请求取消和多rank不同步等边界。判断标准首先是输出与依赖正确，再检查新增同步和额外内存是否抵消减少复制的收益。可更新性与低成本并发需要分别建立证据。

\hypertarget{section-8-3}{%
\subsubsection{8.3 图存档的恢复条件与失效判定}\label{section-8-3}}

Medusa通过分配关联重建直接指针，Foundry通过重现虚拟地址布局和内核环境减少逐项修补。前者受到指针可识别性的约束，后者依赖布局及相应配置能够重建。二者共同表明，保存图拓扑不足以说明一个存档长期有效。模型、算子库、编译选项、KV布局和通信配置变化后，需要知道哪一部分绑定失效，以及失效是否会传播到其他实例\protect\hyperlink{ref-28}{{[}28}, \protect\hyperlink{ref-29}{29{]}}。

值得研究的是按依赖对象进行分层检查：仅改变数据内容时是否可以保留模板，改变内核实现时是否必须重建函数绑定，改变形状上界或通信拓扑时是否需要重新捕获。应在这些变化下同时记录恢复成功率、错误检测能力、回退成本和检查开销。通用进程恢复与图专项物化可以提供不同的状态覆盖范围，但不能仅以恢复时间判断适用性；先明确恢复对象，再比较保留状态与重新准备的成本。

\hypertarget{section-8-4}{%
\subsubsection{8.4 多设备依赖与数据所有权协同}\label{section-8-4}}

Mustard区分跨节点静态分配与节点内动态分配，Foundry的多rank共享依赖等价图结构，MPK则以事件就绪推进更细粒度计算和通信。这些证据覆盖不同的任务分配和状态条件，尚不能由一个``支持多设备''标签概括。LLM中的张量、专家和流水线并行具有不同的图结构与数据驻留要求；尤其是动态任务转移，可能同时引起KV或其他输入的移动\protect\hyperlink{ref-21}{{[}21}, \protect\hyperlink{ref-24}{24}, \protect\hyperlink{ref-29}{29{]}}。

后续比较应固定并行策略，分别考察依赖通知、数据迁移和任务执行位置调整的成本。若运行时重新分配只减少局部等待，却引入更长的数据传输，负载均衡改善不必然缩短迭代。对于模板共享，还应验证各rank的参数和通信状态是否可独立更新，以及不同流水线阶段能否满足结构兼容条件。节点内与跨节点结果需要分别报告。

\hypertarget{section-8-5}{%
\subsubsection{8.5 从稳态加速转向完整使用周期证据}\label{section-8-5}}

表8和表9显示，已有评测分别覆盖内核、模型、任务图和服务启动等窗口。使这些结果形成可用的综述结论，需要清楚区分准备前后、图命中与未命中、单项机制与系统组合。特别是原文表格中的负收益或长尾变化，应和平均加速一起保留；它们能够提示方法依赖的范围，而不是作为异常结果从比较中删除。

共同评测可以从相同模型和请求轨迹出发，分别测量无图、完全预热、按需捕获、恢复后执行及配置变化后的回退。记录主机、设备与通信活动，将图提交、状态更新、并行调度、权重加载和环境预热的贡献拆开；同时报告额外内存、存档大小和输出一致性。这样的证据能够判断一种机制适合长期在线服务、弹性扩容还是低复用短任务，也能帮助确定下一步应补充哪种负载或平台验证。

\hypertarget{section-9}{%
\subsection{9 结束语}\label{section-9}}

面向大语言模型推理的执行图优化涉及图的形成、动态复用、并行执行与实例存续。本文围绕这四个主要决策层次，进一步按具体执行决策归纳十二个小类，并将任务粒度、设备范围、决策时机和验证场景作为横向属性。以机制而非系统名称比较，可以区分哪些方法互为选择、哪些方法能够组合，以及组合时需要满足的状态与资源条件。

已有研究表明，减少逐次提交、缩小更新范围、调整任务分配和迁移准备成本均有明确应用场景，其收益受输入分布、同步关系、内存预算与运行环境制约。原论文的实际评测进一步说明，静态网络加速、Attention吞吐和图恢复时间分别回答不同问题；具备图兼容性或在某一阶段获得加速，都不足以证明完整服务持续受益。

对昇腾而言，xLLM提供了ACLGraph机制与消融证据，版本固定的公开实现补充了分桶、缓存和更新时序的工程约束；CUDA Graph及相邻编译、调度和恢复研究则提供了可逐项比较的设计空间。下一步需要把这些机制所依赖的平台能力、正确性条件和实际成本联系起来，在完整使用周期内判断其适用范围，而非仅以捕获范围或单次回放速度评价进展。

\hypertarget{references}{%
\subsection{参考文献}\label{references}}

\leavevmode\vadjust pre{\hypertarget{ref-1}{}}%
{[}1{]} Gyeong-In Yu, Joo Seong Jeong, Geon-Woo Kim, Soojeong Kim, Byung-Gon Chun. Orca: A Distributed Serving System for Transformer-Based Generative Models. 16th USENIX Symposium on Operating Systems Design and Implementation. 2022. \url{https://www.usenix.org/conference/osdi22/presentation/yu}

\leavevmode\vadjust pre{\hypertarget{ref-2}{}}%
{[}2{]} Woosuk Kwon, Zhuohan Li, Siyuan Zhuang, Ying Sheng, Lianmin Zheng, Cody Hao Yu, Joseph E. Gonzalez, Hao Zhang, Ion Stoica. Efficient Memory Management for Large Language Model Serving with PagedAttention. Proceedings of the 29th Symposium on Operating Systems Principles. 2023. DOI: 10.1145/3600006.3613165. \url{https://arxiv.org/abs/2309.06180v1}

\leavevmode\vadjust pre{\hypertarget{ref-3}{}}%
{[}3{]} NVIDIA. CUDA Programming Guide, 4.2 CUDA Graphs. 官方文档，访问日期2026-11-23. \url{https://docs.nvidia.com/cuda/cuda-programming-guide/04-special-topics/cuda-graphs.html}

\leavevmode\vadjust pre{\hypertarget{ref-4}{}}%
{[}4{]} Huawei. ACL Graph Overview. CANN Community Edition 9.0.0 文档，访问日期2026-11-23. \url{https://www.hiascend.com/doc_center/source/en/CANNCommunityEdition/900/programug/acldevg/runtime_doc_dev_0045.html}

\leavevmode\vadjust pre{\hypertarget{ref-5}{}}%
{[}5{]} Tongxuan Liu, Tao Peng, Peijun Yang, Xiaoyang Zhao, Xiusheng Lu, Weizhe Huang, Zirui Liu, Xiaoyu Chen, Zhiwei Liang, Jun Xiong, Donghe Jin, Minchao Zhang, Jinrong Guo, Yingxu Deng, Xu Zhang, Xianzhe Dong, Siqi Wang, Siyu Wu, Yu Wu, Zihan Tang, Yuting Zeng, Yanshu Wang, Jinguang Liu, Meng Kang, Menxin Li, Yunlong Wang, Yiming Liu, Xiaolong Ma, Yifan Wang, Yichen Zhang, Jinrun Yin, Keyang Zheng, Jiawei Yin, Jun Zhang, Ziyue Wang, Xiaobo Lin, Liangyu Liu, Liwei Lan, Yang Liu, Chunhua Peng, Han Liu, Songcheng Ren, Xuezhu Wang, Yunheng Shen, Yi Wang, Guyue Liu, Yitao Hu, Hui Chen, Tong Yang, Hailong Yang, Jing Li, Guiguang Ding, Ke Zhang. xLLM Technical Report. arXiv:2510.14686v2，2026-03-03；首发于2025年. 2026. DOI: 10.48550/arXiv.2510.14686. \url{https://arxiv.org/abs/2510.14686v2}

\leavevmode\vadjust pre{\hypertarget{ref-6}{}}%
{[}6{]} Abhishek Ghosh, Ajay Nayak, Ashish Panwar, Arkaprava Basu. GraCE: Unlocking CUDA Graphs with Compiler Support for ML Workloads. 20th USENIX Symposium on Operating Systems Design and Implementation. 2026. \url{https://www.usenix.org/conference/osdi26/presentation/ghosh}

\leavevmode\vadjust pre{\hypertarget{ref-7}{}}%
{[}7{]} Zhihang Yuan, Yuzhang Shang, Yang Zhou, Zhen Dong, Zhe Zhou, Chenhao Xue, Bingzhe Wu, Zhikai Li, Qingyi Gu, Yong Jae Lee, Yan Yan, Beidi Chen, Guangyu Sun, Kurt Keutzer. LLM Inference Unveiled: Survey and Roofline Model Insights. arXiv:2402.16363v6. 2024. \url{https://arxiv.org/abs/2402.16363v6}

\leavevmode\vadjust pre{\hypertarget{ref-8}{}}%
{[}8{]} Ranran Zhen, Juntao Li, Yixin Ji, Zhenlin Yang, Tong Liu, Qingrong Xia, Xinyu Duan, Zhefeng Wang, Baoxing Huai, Min Zhang. Taming the Titans: A Survey of Efficient LLM Inference Serving. INLG 2025, pp.522--541. 2025. \url{https://aclanthology.org/2025.inlg-main.32/}

\leavevmode\vadjust pre{\hypertarget{ref-9}{}}%
{[}9{]} Jason Ansel, Edward Yang, Horace He, Natalia Gimelshein, Animesh Jain, Michael Voznesensky, Bin Bao, Peter Bell, David Berard, Evgeni Burovski, Geeta Chauhan, Anjali Chourdia, Will Constable, Alban Desmaison, Zachary DeVito, Elias Ellison, Will Feng, Jiong Gong, Michael Gschwind, Brian Hirsh, Sherlock Huang, Kshiteej Kalambarkar, Laurent Kirsch, Michael Lazos, Mario Lezcano, Yanbo Liang, Jason Liang, Yinghai Lu, C. K. Luk, Bert Maher, Yunjie Pan, Christian Puhrsch, Matthias Reso, Mark Saroufim, Marcos Yukio Siraichi, Helen Suk, Shunting Zhang, Michael Suo, Phil Tillet, Xu Zhao, Eikan Wang, Keren Zhou, Richard Zou, Xiaodong Wang, Ajit Mathews, William Wen, Gregory Chanan, Peng Wu, Soumith Chintala. PyTorch 2: Faster Machine Learning Through Dynamic Python Bytecode Transformation and Graph Compilation. Proceedings of the 29th ACM International Conference on Architectural Support for Programming Languages and Operating Systems, Volume 2. 2024. DOI: 10.1145/3620665.3640366. \url{https://doi.org/10.1145/3620665.3640366}

\leavevmode\vadjust pre{\hypertarget{ref-10}{}}%
{[}10{]} Ruihang Lai, Junru Shao, Siyuan Feng, Steven Lyubomirsky, Bohan Hou, Wuwei Lin, Zihao Ye, Hongyi Jin, Yuchen Jin, Jiawei Liu, Lesheng Jin, Yaxing Cai, Ziheng Jiang, Yong Wu, Sunghyun Park, Prakalp Srivastava, Jared Roesch, Todd C. Mowry, Tianqi Chen. Relax: Composable Abstractions for End-to-End Dynamic Machine Learning. Proceedings of the 30th ACM International Conference on Architectural Support for Programming Languages and Operating Systems, Volume 2. 2025. DOI: 10.1145/3676641.3716249. \url{https://doi.org/10.1145/3676641.3716249}

\leavevmode\vadjust pre{\hypertarget{ref-11}{}}%
{[}11{]} Woosuk Kwon, Gyeong-In Yu, Eunji Jeong, Byung-Gon Chun. Nimble: Lightweight and Parallel GPU Task Scheduling for Deep Learning. Advances in Neural Information Processing Systems 33. 2020. \url{https://proceedings.neurips.cc/paper/2020/hash/5f0ad4db43d8723d18169b2e4817a160-Abstract.html}

\leavevmode\vadjust pre{\hypertarget{ref-12}{}}%
{[}12{]} Tsung-Wei Huang, Dian-Lun Lin, Chun-Xun Lin, Yibo Lin. Taskflow: A Lightweight Parallel and Heterogeneous Task Graph Computing System. IEEE Transactions on Parallel and Distributed Systems 33(6):1303--1320. 2022. DOI: 10.1109/TPDS.2021.3104255. \url{https://doi.org/10.1109/TPDS.2021.3104255}

\leavevmode\vadjust pre{\hypertarget{ref-13}{}}%
{[}13{]} Zihao Ye, Lequn Chen, Ruihang Lai, Wuwei Lin, Yineng Zhang, Stephanie Wang, Tianqi Chen, Baris Kasikci, Vinod Grover, Arvind Krishnamurthy, Luis Ceze. FlashInfer: Efficient and Customizable Attention Engine for LLM Inference Serving. Proceedings of Machine Learning and Systems. 2025. \url{https://proceedings.mlsys.org/paper_files/paper/2025/hash/dbf02b21d77409a2db30e56866a8ab3a-Abstract-Conference.html}

\leavevmode\vadjust pre{\hypertarget{ref-14}{}}%
{[}14{]} Jordan Juravsky, Bradley Brown, Ryan Ehrlich, Daniel Y. Fu, Christopher Ré, Azalia Mirhoseini. Hydragen: High-Throughput LLM Inference with Shared Prefixes. arXiv:2402.05099v2. 2024. DOI: 10.48550/arXiv.2402.05099. \url{https://arxiv.org/abs/2402.05099v2}

\leavevmode\vadjust pre{\hypertarget{ref-15}{}}%
{[}15{]} Haichen Shen, Jared Roesch, Zhi Chen, Wei Chen, Yong Wu, Mu Li, Vin Sharma, Zachary Tatlock, Yida Wang. Nimble: Efficiently Compiling Dynamic Neural Networks for Model Inference. Proceedings of Machine Learning and Systems 3. 2021. \url{https://proceedings.mlsys.org/paper_files/paper/2021/hash/5b47430e24a5a1f9fe21f0e8eb814131-Abstract.html}

\leavevmode\vadjust pre{\hypertarget{ref-16}{}}%
{[}16{]} Zhuoming Chen, Avner May, Ruslan Svirschevski, Yuhsun Huang, Max Ryabinin, Zhihao Jia, Beidi Chen. Sequoia: Scalable and Robust Speculative Decoding. Advances in Neural Information Processing Systems 37. 2024. DOI: 10.52202/079017-4116. \url{https://proceedings.neurips.cc/paper_files/paper/2024/hash/ea1f5f0878d43ff4fb8bf64ef4a2326c-Abstract-Conference.html}

\leavevmode\vadjust pre{\hypertarget{ref-17}{}}%
{[}17{]} Pratik Fegade, Tianqi Chen, Phillip B. Gibbons, Todd C. Mowry. Cortex: A Compiler for Recursive Deep Learning Models. Proceedings of Machine Learning and Systems 3. 2021. \url{https://proceedings.mlsys.org/paper_files/paper/2021/hash/eca986d585a03890a412587a2f5ccb43-Abstract.html}

\leavevmode\vadjust pre{\hypertarget{ref-18}{}}%
{[}18{]} Pratik Fegade, Tianqi Chen, Phillip B. Gibbons, Todd C. Mowry. ACRoBat: Optimizing Auto-batching of Dynamic Deep Learning at Compile Time. Proceedings of Machine Learning and Systems 6. 2024. \url{https://proceedings.mlsys.org/paper_files/paper/2024/hash/096b1019463f34eb241e87cfce8dfe16-Abstract-Conference.html}

\leavevmode\vadjust pre{\hypertarget{ref-19}{}}%
{[}19{]} Aodong Chen, Fei Xu, Li Han, Yuan Dong, Li Chen, Zhi Zhou, Fangming Liu. Opara: Exploiting Operator Parallelism for Expediting DNN Inference on GPUs. IEEE Transactions on Computers. 2025. DOI: 10.1109/TC.2024.3475589. \url{https://doi.org/10.1109/TC.2024.3475589}

\leavevmode\vadjust pre{\hypertarget{ref-20}{}}%
{[}20{]} Yuxuan Zhao, Qi Sun, Zhuolun He, Yang Bai, Bei Yu. AutoGraph: Optimizing DNN Computation Graph for Parallel GPU Kernel Execution. Proceedings of the AAAI Conference on Artificial Intelligence. 2023. DOI: 10.1609/aaai.v37i9.26343. \url{https://doi.org/10.1609/aaai.v37i9.26343}

\leavevmode\vadjust pre{\hypertarget{ref-21}{}}%
{[}21{]} Ilyas Turimbetov, Mohamed Wahib, Didem Unat. A Device-Side Execution Model for Multi-GPU Task Graphs. 2025 International Conference on Supercomputing (ICS 2025). 2025. DOI: 10.1145/3721145.3730426. \url{https://doi.org/10.1145/3721145.3730426}

\leavevmode\vadjust pre{\hypertarget{ref-22}{}}%
{[}22{]} Yaoyao Ding, Ligeng Zhu, Zhihao Jia, Gennady Pekhimenko, Song Han. IOS: Inter-Operator Scheduler for CNN Acceleration. Proceedings of Machine Learning and Systems 3. 2021. \url{https://proceedings.mlsys.org/paper_files/paper/2021/hash/1f8053a67ec8e0b57455713cefdd8218-Abstract.html}

\leavevmode\vadjust pre{\hypertarget{ref-23}{}}%
{[}23{]} Lingxiao Ma, Zhiqiang Xie, Zhi Yang, Jilong Xue, Youshan Miao, Wei Cui, Wenxiang Hu, Fan Yang, Lintao Zhang, Lidong Zhou. Rammer: Enabling Holistic Deep Learning Compiler Optimizations with rTasks. 14th USENIX Symposium on Operating Systems Design and Implementation (OSDI 20). 2020. \url{https://www.usenix.org/conference/osdi20/presentation/ma}

\leavevmode\vadjust pre{\hypertarget{ref-24}{}}%
{[}24{]} Xinhao Cheng, Zhihao Zhang, Yu Zhou, Jianan Ji, Jinchen Jiang, Zepeng Zhao, Ziruo Xiao, Zihao Ye, Yingyi Huang, Ruihang Lai, Hongyi Jin, Bohan Hou, Mengdi Wu, Yixin Dong, Anthony Yip, Zihao Ye, Songting Wang, Wenqin Yang, Xupeng Miao, Tianqi Chen, Zhihao Jia. MPK: A Compiler and Runtime for Mega-Kernelizing Tensor Programs. arXiv:2512.22219v2. 2026. DOI: 10.48550/arXiv.2512.22219. \url{https://arxiv.org/abs/2512.22219v2}

\leavevmode\vadjust pre{\hypertarget{ref-25}{}}%
{[}25{]} Kan Zhu, Yufei Gao, Yilong Zhao, Liangyu Zhao, Gefei Zuo, Yile Gu, Dedong Xie, Tian Tang, Qinyu Xu, Zihao Ye, Keisuke Kamahori, Chien-Yu Lin, Ziren Wang, Stephanie Wang, Arvind Krishnamurthy, Baris Kasikci. NanoFlow: Towards Optimal Large Language Model Serving Throughput. 19th USENIX Symposium on Operating Systems Design and Implementation. 2025. \url{https://www.usenix.org/conference/osdi25/presentation/zhu-kan}

\leavevmode\vadjust pre{\hypertarget{ref-26}{}}%
{[}26{]} Aditya K Kamath, Ramya Prabhu, Jayashree Mohan, Simon Peter, Ramachandran Ramjee, Ashish Panwar. POD-Attention: Unlocking Full Prefill-Decode Overlap for Faster LLM Inference. Proceedings of the 30th ACM International Conference on Architectural Support for Programming Languages and Operating Systems, Volume 2. 2025. DOI: 10.1145/3676641.3715996. \url{https://doi.org/10.1145/3676641.3715996}

\leavevmode\vadjust pre{\hypertarget{ref-27}{}}%
{[}27{]} Divakar Kumar Yadav, Tian Zhao. Hybrid JIT-CUDA Graph Optimization for Low-Latency Large Language Model Inference. arXiv:2604.23467v1. 2026. DOI: 10.48550/arXiv.2604.23467. \url{https://arxiv.org/abs/2604.23467v1}

\leavevmode\vadjust pre{\hypertarget{ref-28}{}}%
{[}28{]} Shaoxun Zeng, Minhui Xie, Shiwei Gao, Youmin Chen, Youyou Lu. Medusa: Accelerating Serverless LLM Inference with Materialization. Proceedings of ASPLOS 2025, Volume 1. 2025. DOI: 10.1145/3669940.3707285. \url{https://doi.org/10.1145/3669940.3707285}

\leavevmode\vadjust pre{\hypertarget{ref-29}{}}%
{[}29{]} Xueshen Liu, Yongji Wu, Yuncheng Yao, Danyang Zhuo, Ion Stoica, Z. Morley Mao. Foundry: Template-Based CUDA Graph Context Materialization for Fast LLM Serving Cold Start. arXiv:2604.06664v1. 2026. DOI: 10.48550/arXiv.2604.06664. \url{https://arxiv.org/abs/2604.06664v1}

\leavevmode\vadjust pre{\hypertarget{ref-30}{}}%
{[}30{]} Junhao Hu, Jiang Xu, Zhixia Liu, Yulong He, Yuetao Chen, Hao Xu, Jiang Liu, Jie Meng, Baoquan Zhang, Shining Wan, Gengyuan Dan, Zhiyu Dong, Zhihao Ren, Changhong Liu, Tao Xie, Dayun Lin, Qin Zhang, Yue Yu, Hao Feng, Xusheng Chen, Yizhou Shan. DEEPSERVE: Serverless Large Language Model Serving at Scale. 2025 USENIX Annual Technical Conference (USENIX ATC 25). 2025. \url{https://www.usenix.org/conference/atc25/presentation/hu-junhao}

\leavevmode\vadjust pre{\hypertarget{ref-31}{}}%
{[}31{]} Xingda Wei, Zhuobin Huang, Tianle Sun, Yingyi Hao, Rong Chen, Mingcong Han, Jinyu Gu, Haibo Chen. PhoenixOS: Concurrent OS-level GPU Checkpoint and Restore with Validated Speculation. Proceedings of the ACM SIGOPS 31st Symposium on Operating Systems Principles (SOSP 25). 2025. DOI: 10.1145/3731569.3764813. \url{https://doi.org/10.1145/3731569.3764813}

\leavevmode\vadjust pre{\hypertarget{ref-32}{}}%
{[}32{]} Ying Sheng, Shiyi Cao, Dacheng Li, Coleman Hooper, Nicholas Lee, Shuo Yang, Christopher Chou, Banghua Zhu, Lianmin Zheng, Kurt Keutzer, Joseph E. Gonzalez, Ion Stoica. S-LoRA: Serving Thousands of Concurrent LoRA Adapters. MLSys 2024. 2024. \url{https://proceedings.mlsys.org/paper_files/paper/2024/hash/906419cd502575b617cc489a1a696a67-Abstract-Conference.html}

\leavevmode\vadjust pre{\hypertarget{ref-33}{}}%
{[}33{]} Lequn Chen, Zihao Ye, Yongji Wu, Danyang Zhuo, Luis Ceze, Arvind Krishnamurthy. Punica: Multi-Tenant LoRA Serving. MLSys 2024. 2024. \url{https://proceedings.mlsys.org/paper_files/paper/2024/hash/054de805fcceb78a201f5e9d53c85908-Abstract-Conference.html}

\leavevmode\vadjust pre{\hypertarget{ref-34}{}}%
{[}34{]} Amey Agrawal, Nitin Kedia, Ashish Panwar, Jayashree Mohan, Nipun Kwatra, Bhargav Gulavani, Alexey Tumanov, Ramachandran Ramjee. Taming Throughput-Latency Tradeoff in LLM Inference with Sarathi-Serve. 18th USENIX Symposium on Operating Systems Design and Implementation. 2024. \url{https://www.usenix.org/conference/osdi24/presentation/agrawal}

\leavevmode\vadjust pre{\hypertarget{ref-35}{}}%
{[}35{]} Kai Zhu, Wenyi Zhao, Zhen Zheng, Tianyou Guo, Pengzhan Zhao, Feiwen Zhu, Junjie Bai, Jun Yang, Xiaoyong Liu, Lansong Diao, Wei Lin. DISC: A Dynamic Shape Compiler for Machine Learning Workloads. Proceedings of the 1st Workshop on Machine Learning and Systems. 2021. DOI: 10.1145/3437984.3458888. \url{https://doi.org/10.1145/3437984.3458888}

\leavevmode\vadjust pre{\hypertarget{ref-36}{}}%
{[}36{]} Pratik Fegade, Tianqi Chen, Phillip B. Gibbons, Todd C. Mowry. The CoRa Tensor Compiler: Compilation for Ragged Tensors with Minimal Padding. Proceedings of Machine Learning and Systems 4. 2022. \url{https://www.pdl.cmu.edu/PDL-FTP/BigLearning/MLSys-2022-the-cora-tensor-compiler.pdf}

\leavevmode\vadjust pre{\hypertarget{ref-37}{}}%
{[}37{]} Xupeng Miao, Gabriele Oliaro, Zhihao Zhang, Xinhao Cheng, Zeyu Wang, Zhengxin Zhang, Rae Ying Yee Wong, Alan Zhu, Lijie Yang, Xiaoxiang Shi, Chunan Shi, Zhuoming Chen, Daiyaan Arfeen, Reyna Abhyankar, Zhihao Jia. SpecInfer: Accelerating Large Language Model Serving with Tree-based Speculative Inference and Verification. Proceedings of ASPLOS 2024, Volume 3. 2024. DOI: 10.1145/3620666.3651335. \url{https://arxiv.org/abs/2305.09781v4}

\leavevmode\vadjust pre{\hypertarget{ref-38}{}}%
{[}38{]} Yuhui Li, Fangyun Wei, Chao Zhang, Hongyang Zhang. EAGLE-2: Faster Inference of Language Models with Dynamic Draft Trees. Proceedings of EMNLP 2024. 2024. DOI: 10.18653/v1/2024.emnlp-main.422. \url{https://aclanthology.org/2024.emnlp-main.422/}

\leavevmode\vadjust pre{\hypertarget{ref-39}{}}%
{[}39{]} Ramya Prabhu, Ajay Nayak, Jayashree Mohan, Ramchandran Ramjee, Ashish Panwar. vAttention: Dynamic Memory Management for Serving LLMs without PagedAttention. Proceedings of ASPLOS 2025, Volume 1. 2025. DOI: 10.1145/3669940.3707256. \url{https://arxiv.org/abs/2405.04437v3}

\leavevmode\vadjust pre{\hypertarget{ref-40}{}}%
{[}40{]} Mengdi Wu, Xinhao Cheng, Shengyu Liu, Chunan Shi, Jianan Ji, Man Kit Ao, Praveen Velliengiri, Xupeng Miao, Oded Padon, Zhihao Jia. Mirage: A Multi-Level Superoptimizer for Tensor Programs. 19th USENIX Symposium on Operating Systems Design and Implementation (OSDI 25). 2025. \url{https://www.usenix.org/conference/osdi25/presentation/wu-mengdi}

\leavevmode\vadjust pre{\hypertarget{ref-41}{}}%
{[}41{]} Radostin Stoyanov, Viktória Spišaková, Jesus Ramos, Steven Gurfinkel, Andrei Vagin, Adrian Reber, Wesley Armour, Rodrigo Bruno. CRIUgpu: Transparent Checkpointing of GPU-Accelerated Workloads. arXiv:2502.16631v1. 2025. DOI: 10.48550/arXiv.2502.16631. \url{https://arxiv.org/abs/2502.16631v1}

\leavevmode\vadjust pre{\hypertarget{ref-42}{}}%
{[}42{]} Shaoxun Zeng, Tingxu Ren, Jiwu Shu, Youyou Lu. GPU Checkpoint/Restore Made Fast and Lightweight. 24th USENIX Conference on File and Storage Technologies (FAST 26). 2026. \url{https://www.usenix.org/conference/fast26/presentation/zeng}

\leavevmode\vadjust pre{\hypertarget{ref-43}{}}%
{[}43{]} Ascend/pytorch Contributors. {[}feat{]} aclGraph task group. 已合并 PR !20264，v2.4.0 分支；访问日期2026-11-23. 2025. \url{https://gitee.com/ascend/pytorch/pulls/20264}

\leavevmode\vadjust pre{\hypertarget{ref-44}{}}%
{[}44{]} vLLM-Ascend Team. ACL Graph. v0.13.0 官方文档，访问日期2026-11-23. \url{https://docs.vllm.ai/projects/ascend/en/v0.13.0/developer_guide/feature_guide/ACL_Graph.html}

\leavevmode\vadjust pre{\hypertarget{ref-45}{}}%
{[}45{]} vLLM-Ascend Contributors. ACLGraphWrapper and attention parameter update implementation. v0.13.0，提交6281c1207a7a499e9f23a42b3a1e7027469f2b10；访问日期2026-11-23. \url{https://github.com/vllm-project/vllm-ascend/blob/6281c1207a7a499e9f23a42b3a1e7027469f2b10/vllm_ascend/compilation/acl_graph.py}

\leavevmode\vadjust pre{\hypertarget{ref-46}{}}%
{[}46{]} Euijun Chung, Yuxiao Jia, Aaron Jezghani, Hyesoon Kim. Characterizing CPU-Induced Slowdowns in Multi-GPU LLM Inference. IISWC 2026 accepted camera-ready; arXiv v3. 2026. \url{https://arxiv.org/abs/2603.22774v3}

\leavevmode\vadjust pre{\hypertarget{ref-47}{}}%
{[}47{]} David Pape, Jonathan Evertz, Lea Schönherr. The Silent Hyperparameter: Quantifying the Impact of Inference Backends on LLM Reproducibility. arXiv:2605.19537v4. 2026. \url{https://arxiv.org/abs/2605.19537v4}

\end{document}
